\documentclass[11pt]{article}
\usepackage[margin=1in]{geometry}
\usepackage{amsmath,amssymb,booktabs,graphicx,xcolor,hyperref,natbib,microtype}
\usepackage{times}
\hypersetup{colorlinks=true,linkcolor=blue!60!black,citecolor=blue!60!black,urlcolor=blue!60!black}
\newcommand{\pal}{\textsc{Palimpsest}}
\newcommand{\rs}{\textsc{Revise-Stream}}
\newcommand{\todo}[1]{\textcolor{red}{[#1]}}
\newcommand{\sOnePropEffect}{20.6}
\newcommand{\sOnePropLo}{18.7}
\newcommand{\sOnePropHi}{22.5}
\newcommand{\sOnePropN}{495}
\newcommand{\sOneCommitUnsup}{9.0}
\newcommand{\sOneCommitUnsupLo}{7.8}
\newcommand{\sOneCommitUnsupHi}{10.2}
\newcommand{\sOneCommitRate}{78}
\newcommand{\sOneAbstainRes}{0.000}
\newcommand{\sOneAbstainUnres}{0.393}
\newcommand{\sOneCommitRes}{0.770}
\newcommand{\sOneCommitUnres}{0.000}
\newcommand{\sOneTdbWarr}{19.3}
\newcommand{\sOneTdbCommit}{85}
\newcommand{\sOneProvP}{0.99}
\newcommand{\sOneProvR}{0.99}
\newcommand{\sOneWarrPerStream}{7.4}
\newcommand{\sOneAnswers}{16683}
\newcommand{\sOneLitPZeroEffect}{36.2}
\newcommand{\sOneLitPZeroWarr}{3.4}
\newcommand{\sOnePOneEffect}{31.2}
\newcommand{\sOnePOneWarr}{3.5}
\newcommand{\RqLearnedEce}{0.113}
\newcommand{\RqLearnedBrier}{0.175}
\newcommand{\RqLearnedAurc}{0.131}
\newcommand{\RqPriorEce}{0.227}
\newcommand{\RqPriorBrier}{0.255}
\newcommand{\RqPriorAurc}{0.168}
\newcommand{\RqLogregEce}{0.034}
\newcommand{\RqLogregBrier}{0.138}
\newcommand{\RqLogregAurc}{0.112}
\newcommand{\RqDefaultEce}{0.118}
\newcommand{\RqDefaultBrier}{0.182}
\newcommand{\RqDefaultAurc}{0.138}
\newcommand{\RqNoRecEce}{0.202}
\newcommand{\RqNoRecBrier}{0.236}
\newcommand{\RqNoRecAurc}{0.156}
\newcommand{\RqNoOrigEce}{0.121}
\newcommand{\RqNoOrigBrier}{0.179}
\newcommand{\RqNoOrigAurc}{0.136}
\newcommand{\RqN}{2481}
\newcommand{\PoisHarmPal}{11.2}
\newcommand{\PoisHarmTdb}{15.1}
\newcommand{\PoisHarmCommit}{13.7}
\newcommand{\PoisSupLowPois}{0.21--0.33}
\newcommand{\PoisSupLowHon}{0.65--0.67}
\newcommand{\PoisSupTrPois}{0.63--0.66}
\newcommand{\PoisSupTrHon}{0.63--0.67}
\newcommand{\sTwoHOne}{44.2}
\newcommand{\sTwoHOneLo}{39.0}
\newcommand{\sTwoHOneHi}{49.7}
\newcommand{\sTwoHOneN}{99}
\newcommand{\sTwoProp}{20.7}
\newcommand{\sTwoPropLo}{16.2}
\newcommand{\sTwoPropHi}{25.4}
\newcommand{\sTwoPalWarr}{0.956}
\newcommand{\sTwoPalRes}{0.976}
\newcommand{\sTwoPalUnres}{0.834}
\newcommand{\sTwoPalCommit}{0.080}
\newcommand{\sTwoPalUnsup}{0.083}
\newcommand{\sTwoPalCost}{0.078}
\newcommand{\sTwoPalTok}{188}
\newcommand{\sTwoLogWarr}{0.514}
\newcommand{\sTwoLogRes}{0.666}
\newcommand{\sTwoLogUnres}{0.128}
\newcommand{\sTwoLogCommit}{0.572}
\newcommand{\sTwoLogUnsup}{0.392}
\newcommand{\sTwoLogCost}{0.234}
\newcommand{\sTwoLogTok}{532}
\newcommand{\sTwoNlWarr}{0.330}
\newcommand{\sTwoNlRes}{0.554}
\newcommand{\sTwoNlUnres}{0.101}
\newcommand{\sTwoNlCommit}{0.606}
\newcommand{\sTwoNlUnsup}{0.456}
\newcommand{\sTwoNlCost}{0.203}
\newcommand{\sTwoNlTok}{439}
\newcommand{\sTwoTdbWarr}{0.779}
\newcommand{\sTwoTdbRes}{0.895}
\newcommand{\sTwoTdbUnres}{0.000}
\newcommand{\sTwoTdbCommit}{0.837}
\newcommand{\sTwoTdbUnsup}{0.142}
\newcommand{\sTwoTdbCost}{0.078}
\newcommand{\sTwoTdbTok}{188}
\newcommand{\sTwoCommitWarr}{0.979}
\newcommand{\sTwoCommitRes}{0.987}
\newcommand{\sTwoCommitUnres}{0.045}
\newcommand{\sTwoCommitCommit}{0.758}
\newcommand{\sTwoCommitUnsup}{0.104}
\newcommand{\sTwoCommitCost}{0.078}
\newcommand{\sTwoCommitTok}{188}
\newcommand{\sTwoTopkWarr}{0.419}
\newcommand{\sTwoTopkRes}{0.614}
\newcommand{\sTwoTopkUnres}{0.103}
\newcommand{\sTwoTopkCommit}{0.504}
\newcommand{\sTwoTopkUnsup}{0.533}
\newcommand{\sTwoTopkCost}{0.196}
\newcommand{\sTwoTopkTok}{454}
\newcommand{\sTwoNopropWarr}{0.749}
\newcommand{\sTwoNopropRes}{0.920}
\newcommand{\sTwoNopropUnres}{0.813}
\newcommand{\sTwoNopropCommit}{0.079}
\newcommand{\sTwoNopropUnsup}{0.093}
\newcommand{\sTwoNopropCost}{0.078}
\newcommand{\sTwoNopropTok}{188}
\newcommand{\sTwoRuleWarr}{0.956}
\newcommand{\sTwoRuleRes}{0.976}
\newcommand{\sTwoRuleUnres}{0.834}
\newcommand{\sTwoRuleCommit}{0.080}
\newcommand{\sTwoRuleUnsup}{0.083}
\newcommand{\sTwoRuleCost}{0.078}
\newcommand{\sTwoRuleTok}{188}
\newcommand{\sTwoExtSlot}{1.5}
\newcommand{\sTwoExtMissed}{3.7}
\newcommand{\sTwoExtSpur}{0.0}
\newcommand{\sTwoExtValid}{6.0}
\newcommand{\sTwoExtFail}{0.1}
\newcommand{\sTwoWarrPerStream}{7.41}
\newcommand{\sTwoSubPalWarr}{0.989}
\newcommand{\sTwoSubN}{30}
\newcommand{\sTwoCtwoLogWarr}{0.400}
\newcommand{\sTwoCtwoLogDirect}{0.703}
\newcommand{\sTwoCtwoLogRes}{0.607}
\newcommand{\sTwoCtwoLogUnres}{0.134}
\newcommand{\sTwoCtwoLogCommit}{0.533}
\newcommand{\sTwoCtwoLogMalformed}{0.000}
\newcommand{\sTwoCtwoLogTok}{552}
\newcommand{\sTwoCtwoLogCost}{0.243}
\newcommand{\sTwoCtwoLogDiff}{59.0}
\newcommand{\sTwoCtwoLogDiffLo}{51.4}
\newcommand{\sTwoCtwoLogDiffHi}{66.5}
\newcommand{\sTwoCtwoLogDiffN}{30}
\newcommand{\sTwoCtwoTopTenWarr}{0.225}
\newcommand{\sTwoCtwoTopTenDirect}{0.787}
\newcommand{\sTwoCtwoTopTenRes}{0.572}
\newcommand{\sTwoCtwoTopTenUnres}{0.138}
\newcommand{\sTwoCtwoTopTenCommit}{0.377}
\newcommand{\sTwoCtwoTopTenMalformed}{0.000}
\newcommand{\sTwoCtwoTopTenTok}{399}
\newcommand{\sTwoCtwoTopTenCost}{0.171}
\newcommand{\sTwoCtwoTopTenDiff}{76.5}
\newcommand{\sTwoCtwoTopTenDiffLo}{69.6}
\newcommand{\sTwoCtwoTopTenDiffHi}{82.7}
\newcommand{\sTwoCtwoTopTenDiffN}{30}
\newcommand{\sTwoCtwoTopTwentyWarr}{0.385}
\newcommand{\sTwoCtwoTopTwentyDirect}{0.723}
\newcommand{\sTwoCtwoTopTwentyRes}{0.610}
\newcommand{\sTwoCtwoTopTwentyUnres}{0.136}
\newcommand{\sTwoCtwoTopTwentyCommit}{0.412}
\newcommand{\sTwoCtwoTopTwentyMalformed}{0.000}
\newcommand{\sTwoCtwoTopTwentyTok}{424}
\newcommand{\sTwoCtwoTopTwentyCost}{0.182}
\newcommand{\sTwoCtwoTopTwentyDiff}{60.4}
\newcommand{\sTwoCtwoTopTwentyDiffLo}{51.4}
\newcommand{\sTwoCtwoTopTwentyDiffHi}{69.3}
\newcommand{\sTwoCtwoTopTwentyDiffN}{30}
\newcommand{\sTwoCtwoTopFortyWarr}{0.429}
\newcommand{\sTwoCtwoTopFortyDirect}{0.739}
\newcommand{\sTwoCtwoTopFortyRes}{0.628}
\newcommand{\sTwoCtwoTopFortyUnres}{0.133}
\newcommand{\sTwoCtwoTopFortyCommit}{0.471}
\newcommand{\sTwoCtwoTopFortyMalformed}{0.000}
\newcommand{\sTwoCtwoTopFortyTok}{464}
\newcommand{\sTwoCtwoTopFortyCost}{0.201}
\newcommand{\sTwoCtwoTopFortyDiff}{56.0}
\newcommand{\sTwoCtwoTopFortyDiffLo}{45.7}
\newcommand{\sTwoCtwoTopFortyDiffHi}{66.2}
\newcommand{\sTwoCtwoTopFortyDiffN}{30}
\newcommand{\sTwoCtwoNlWarr}{0.323}
\newcommand{\sTwoCtwoNlDirect}{0.672}
\newcommand{\sTwoCtwoNlRes}{0.524}
\newcommand{\sTwoCtwoNlUnres}{0.102}
\newcommand{\sTwoCtwoNlCommit}{0.563}
\newcommand{\sTwoCtwoNlMalformed}{0.000}
\newcommand{\sTwoCtwoNlTok}{448}
\newcommand{\sTwoCtwoNlCost}{0.207}
\newcommand{\sTwoCtwoNlDiff}{66.7}
\newcommand{\sTwoCtwoNlDiffLo}{56.5}
\newcommand{\sTwoCtwoNlDiffHi}{76.2}
\newcommand{\sTwoCtwoNlDiffN}{30}
\newcommand{\sTwoSubTopFortyWarr}{0.419}
\newcommand{\sTwoSubTopFortyDirect}{0.731}
\newcommand{\sTwoSubTopFortyRes}{0.614}
\newcommand{\sTwoSubTopFortyUnres}{0.103}
\newcommand{\sTwoSubTopFortyCommit}{0.504}
\newcommand{\sTwoSubTopFortyMalformed}{0.000}
\newcommand{\sTwoSubTopFortyTok}{454}
\newcommand{\sTwoSubTopFortyCost}{0.196}
\newcommand{\sTwoSubTopFortyDiff}{57.0}
\newcommand{\sTwoSubTopFortyDiffLo}{46.8}
\newcommand{\sTwoSubTopFortyDiffHi}{67.1}
\newcommand{\sTwoSubTopFortyDiffN}{30}
\newcommand{\sTwoSubNlWarr}{0.330}
\newcommand{\sTwoSubNlDirect}{0.685}
\newcommand{\sTwoSubNlRes}{0.554}
\newcommand{\sTwoSubNlUnres}{0.101}
\newcommand{\sTwoSubNlCommit}{0.606}
\newcommand{\sTwoSubNlMalformed}{0.000}
\newcommand{\sTwoSubNlTok}{439}
\newcommand{\sTwoSubNlCost}{0.203}
\newcommand{\sTwoSubNlDiff}{66.0}
\newcommand{\sTwoSubNlDiffLo}{55.8}
\newcommand{\sTwoSubNlDiffHi}{75.3}
\newcommand{\sTwoSubNlDiffN}{30}
\newcommand{\sTwoSubLogDiff}{58.2}
\newcommand{\sTwoSubLogDiffLo}{50.5}
\newcommand{\sTwoSubLogDiffHi}{66.3}
\newcommand{\sTwoSubLogWarr}{0.408}
\newcommand{\sThreeCtwoLogWarr}{0.019}
\newcommand{\sThreeCtwoLogDirect}{0.491}
\newcommand{\sThreeCtwoLogRes}{0.205}
\newcommand{\sThreeCtwoLogUnres}{0.394}
\newcommand{\sThreeCtwoLogCommit}{0.485}
\newcommand{\sThreeCtwoNlWarr}{0.173}
\newcommand{\sThreeCtwoNlDirect}{0.466}
\newcommand{\sThreeCtwoNlRes}{0.208}
\newcommand{\sThreeCtwoNlUnres}{0.318}
\newcommand{\sThreeCtwoNlCommit}{0.439}
\newcommand{\sStrongHOne}{18.7}
\newcommand{\sStrongHOneLo}{7.4}
\newcommand{\sStrongHOneHi}{33.1}
\newcommand{\sStrongPalWarr}{1.000}
\newcommand{\sStrongLogWarr}{0.813}
\newcommand{\sStrongLogCommit}{0.489}
\newcommand{\sStrongExtSlot}{0.3}
\newcommand{\xE}{188}
\newcommand{\xAnl}{6.5}
\newcommand{\xAlog}{5.1}
\newcommand{\xRstar}{0.22}
\newcommand{\xNobs}{131}
\newcommand{\RqTwoSupEce}{0.150}
\newcommand{\RqTwoSupBrier}{0.133}
\newcommand{\RqTwoSupAurc}{0.073}
\newcommand{\RqTwoSupN}{973}
\newcommand{\RqTwoVerbEce}{0.088}
\newcommand{\RqTwoVerbBrier}{0.091}
\newcommand{\RqTwoVerbAurc}{0.069}
\newcommand{\RqTwoVerbN}{973}
\newcommand{\RqTwoSampEce}{0.073}
\newcommand{\RqTwoSampBrier}{0.083}
\newcommand{\RqTwoSampAurc}{0.078}
\newcommand{\RqTwoSampN}{300}
\newcommand{\nzSlotTwentyPal}{0.449}
\newcommand{\nzSlotTwentyTdb}{0.388}
\newcommand{\nzSlotThirtyfivePal}{0.273}
\newcommand{\nzSlotThirtyfiveTdb}{0.301}
\newcommand{\nzCueFiftyPal}{0.562}
\newcommand{\nzCueFiftyTdb}{0.758}
\newcommand{\sThreeHOne}{90.4}
\newcommand{\sThreeHOneLo}{73.1}
\newcommand{\sThreeHOneHi}{100.0}
\newcommand{\sThreeHOneN}{13}
\newcommand{\sThreeProp}{19.9}
\newcommand{\sThreePropLo}{2.6}
\newcommand{\sThreePropHi}{41.0}
\newcommand{\sThreeDirect}{45.7}
\newcommand{\sThreeDirectLo}{31.1}
\newcommand{\sThreeDirectHi}{60.5}
\newcommand{\sThreeDirectN}{27}
\newcommand{\sThreePalWarr}{0.923}
\newcommand{\sThreePalDirect}{0.976}
\newcommand{\sThreePalRes}{0.978}
\newcommand{\sThreePalUnres}{0.909}
\newcommand{\sThreePalCommit}{0.000}
\newcommand{\sThreePalUnsup}{0.107}
\newcommand{\sThreeLogWarr}{0.019}
\newcommand{\sThreeLogDirect}{0.519}
\newcommand{\sThreeLogRes}{0.185}
\newcommand{\sThreeLogUnres}{0.333}
\newcommand{\sThreeLogCommit}{0.545}
\newcommand{\sThreeLogUnsup}{0.476}
\newcommand{\sThreeTdbWarr}{0.865}
\newcommand{\sThreeTdbDirect}{0.884}
\newcommand{\sThreeTdbRes}{0.692}
\newcommand{\sThreeTdbUnres}{0.000}
\newcommand{\sThreeTdbCommit}{1.000}
\newcommand{\sThreeTdbUnsup}{0.107}
\newcommand{\sThreeRuleWarr}{0.923}
\newcommand{\sThreeRuleDirect}{0.976}
\newcommand{\sThreeRuleRes}{0.978}
\newcommand{\sThreeRuleUnres}{0.909}
\newcommand{\sThreeRuleCommit}{0.000}
\newcommand{\sThreeRuleUnsup}{0.107}
\newcommand{\sThreeCommitWarr}{0.923}
\newcommand{\sThreeCommitDirect}{0.976}
\newcommand{\sThreeCommitRes}{0.978}
\newcommand{\sThreeCommitUnres}{0.000}
\newcommand{\sThreeCommitCommit}{1.000}
\newcommand{\sThreeCommitUnsup}{0.107}
\newcommand{\sThreeNlWarr}{0.077}
\newcommand{\sThreeNlDirect}{0.451}
\newcommand{\sThreeNlRes}{0.181}
\newcommand{\sThreeNlUnres}{0.182}
\newcommand{\sThreeNlCommit}{0.576}
\newcommand{\sThreeNlUnsup}{0.446}
\newcommand{\sThreeExtSlot}{1.2}
\newcommand{\sThreeExtMissed}{0.0}
\newcommand{\sThreeExtSpur}{0.0}
\newcommand{\sThreeExtFail}{0.1}
\newcommand{\sThreeWarrPerStream}{0.80}
\newcommand{\sThreeEstTotal}{338}
\newcommand{\sThreeEstEmpty}{94}
\newcommand{\sThreeHoldsFalse}{102}
\newcommand{\sThreeWarrEmpty}{3}
\newcommand{\sThreeWarrTotal}{24}
\newcommand{\sThreePalResNe}{0.954}
\newcommand{\sThreePalWarrNe}{0.923}
\newcommand{\sThreeLogResNe}{0.518}
\newcommand{\sThreeLogWarrNe}{0.077}
\newcommand{\sThreeTdbResNe}{0.954}
\newcommand{\sThreeTdbWarrNe}{0.923}
\newcommand{\lmePalCov}{0.32}
\newcommand{\lmePalAcc}{0.167}
\newcommand{\lmePalAccAns}{0.400}
\newcommand{\lmePalSuCov}{0.74}
\newcommand{\lmePalSuAcc}{0.474}
\newcommand{\lmePalSuAccAns}{0.621}
\newcommand{\lmeTdbCov}{0.88}
\newcommand{\lmeTdbAcc}{0.526}
\newcommand{\lmeTdbAccAns}{0.594}
\newcommand{\lmeLogCov}{0.56}
\newcommand{\lmeLogAcc}{0.359}
\newcommand{\lmeLogAccAns}{0.614}
\newcommand{\lmeCommitCov}{0.88}
\newcommand{\lmeCommitAcc}{0.526}
\newcommand{\lmeCommitAccAns}{0.594}
\newcommand{\lmeN}{78}
\newcommand{\lmeMapped}{69}
\newcommand{\lmePalUnres}{44}
\newcommand{\lmePalEst}{25}
\newcommand{\sysRpOneIncCpu}{0.33}
\newcommand{\sysRpOneRepCpu}{0.03}
\newcommand{\sysRpOneIncLat}{0.4}
\newcommand{\sysRpOneRepLat}{2.6}
\newcommand{\sysRpOneIncPNN}{1.0}
\newcommand{\sysRpOneRepPNN}{6.7}
\newcommand{\sysRpOneIncMax}{12}
\newcommand{\sysRpOneRepMax}{61}
\newcommand{\sysRpOneIncMem}{1.0}
\newcommand{\sysRpOneRepMem}{0.1}
\newcommand{\sysRpThirdIncCpu}{0.32}
\newcommand{\sysRpThirdRepCpu}{0.11}
\newcommand{\sysRpThirdIncLat}{0.4}
\newcommand{\sysRpThirdRepLat}{2.6}
\newcommand{\sysRpThirdIncPNN}{1.2}
\newcommand{\sysRpThirdRepPNN}{9.2}
\newcommand{\sysRpThirdIncMax}{10}
\newcommand{\sysRpThirdRepMax}{69}
\newcommand{\sysRpThirdIncMem}{0.9}
\newcommand{\sysRpThirdRepMem}{0.3}
\newcommand{\sysROneIncCpu}{0.36}
\newcommand{\sysROneRepCpu}{0.32}
\newcommand{\sysROneIncLat}{0.4}
\newcommand{\sysROneRepLat}{2.4}
\newcommand{\sysROneIncPNN}{1.5}
\newcommand{\sysROneRepPNN}{11.0}
\newcommand{\sysROneIncMax}{57}
\newcommand{\sysROneRepMax}{79}
\newcommand{\sysROneIncMem}{1.0}
\newcommand{\sysROneRepMem}{0.5}
\newcommand{\sysRThreeIncCpu}{0.46}
\newcommand{\sysRThreeRepCpu}{0.93}
\newcommand{\sysRThreeIncLat}{0.3}
\newcommand{\sysRThreeRepLat}{2.3}
\newcommand{\sysRThreeIncPNN}{1.6}
\newcommand{\sysRThreeRepPNN}{11.7}
\newcommand{\sysRThreeIncMax}{40}
\newcommand{\sysRThreeRepMax}{85}
\newcommand{\sysRThreeIncMem}{1.2}
\newcommand{\sysRThreeRepMem}{0.7}
\newcommand{\sysMThirteenIncCpu}{3.72}
\newcommand{\sysMThirteenRepCpu}{1.66}
\newcommand{\sysMThirteenIncLat}{0.9}
\newcommand{\sysMThirteenRepLat}{12.5}
\newcommand{\sysMThirteenIncPNN}{1.8}
\newcommand{\sysMThirteenRepPNN}{34.3}
\newcommand{\sysMThirteenIncMax}{215}
\newcommand{\sysMThirteenRepMax}{2553}
\newcommand{\sysMThirteenIncMem}{22.4}
\newcommand{\sysMThirteenRepMem}{4.2}
\newcommand{\sysTTwoFortyIncCpu}{0.79}
\newcommand{\sysTTwoFortyRepCpu}{1.38}
\newcommand{\sysTTwoFortyIncLat}{0.9}
\newcommand{\sysTTwoFortyRepLat}{7.4}
\newcommand{\sysTTwoFortyIncPNN}{4.1}
\newcommand{\sysTTwoFortyRepPNN}{31.7}
\newcommand{\sysTTwoFortyIncMax}{16}
\newcommand{\sysTTwoFortyRepMax}{108}
\newcommand{\sysTTwoFortyIncMem}{2.9}
\newcommand{\sysTTwoFortyRepMem}{1.1}
\newcommand{\sysAnswers}{40030}
\newcommand{\sysDisagree}{0}
\newcommand{\rgTotal}{3789}
\newcommand{\rgLatestTruePct}{62}
\newcommand{\rgLatestWrongPct}{10}
\newcommand{\rgSelfN}{53}
\newcommand{\rgSelfLww}{0.96}
\newcommand{\rgSelfPal}{0.68}
\newcommand{\rgSelfAbst}{0.32}
\newcommand{\rgSelfWrong}{0.00}
\newcommand{\rgSelfArgmax}{0.94}
\newcommand{\rgMultiN}{854}
\newcommand{\rgMultiLww}{0.92}
\newcommand{\rgMultiPal}{0.71}
\newcommand{\rgMultiAbst}{0.24}
\newcommand{\rgMultiWrong}{0.04}
\newcommand{\rgMultiArgmax}{0.85}
\newcommand{\rgStaleN}{1427}
\newcommand{\rgStaleLww}{0.92}
\newcommand{\rgStalePal}{0.56}
\newcommand{\rgStaleAbst}{0.41}
\newcommand{\rgStaleWrong}{0.04}
\newcommand{\rgStaleArgmax}{0.84}
\newcommand{\rgDelayN}{75}
\newcommand{\rgDelayLww}{0.12}
\newcommand{\rgDelayPal}{0.11}
\newcommand{\rgDelayAbst}{0.63}
\newcommand{\rgDelayWrong}{0.27}
\newcommand{\rgDelayArgmax}{0.35}
\newcommand{\rgLowN}{192}
\newcommand{\rgLowLww}{0.03}
\newcommand{\rgLowPal}{0.01}
\newcommand{\rgLowAbst}{0.95}
\newcommand{\rgLowWrong}{0.04}
\newcommand{\rgLowArgmax}{0.69}
\newcommand{\rgOthN}{96}
\newcommand{\rgOthLww}{0.05}
\newcommand{\rgOthPal}{0.01}
\newcommand{\rgOthAbst}{0.91}
\newcommand{\rgOthWrong}{0.08}
\newcommand{\rgOthArgmax}{0.39}
\newcommand{\rgRetrN}{242}
\newcommand{\rgRetrLww}{0.51}
\newcommand{\rgRetrPal}{0.35}
\newcommand{\rgRetrAbst}{0.31}
\newcommand{\rgRetrWrong}{0.34}
\newcommand{\rgRetrArgmax}{0.47}
\newcommand{\rgDisputedN}{1191}
\newcommand{\rgArgmaxFull}{0.635}
\newcommand{\rgLwwFull}{0.630}
\newcommand{\rgHalfCov}{0.51}
\newcommand{\rgHalfArgmax}{0.718}
\newcommand{\rgHalfLww}{0.652}
\newcommand{\rgOneTotal}{6765}
\newcommand{\rgOneLatestTruePct}{58}
\newcommand{\rgOneLatestWrongPct}{11}
\newcommand{\rgOneSelfN}{50}
\newcommand{\rgOneSelfLww}{0.92}
\newcommand{\rgOneSelfPal}{0.78}
\newcommand{\rgOneSelfAbst}{0.20}
\newcommand{\rgOneSelfWrong}{0.02}
\newcommand{\rgOneSelfArgmax}{0.94}
\newcommand{\rgOneMultiN}{1504}
\newcommand{\rgOneMultiLww}{0.92}
\newcommand{\rgOneMultiPal}{0.69}
\newcommand{\rgOneMultiAbst}{0.27}
\newcommand{\rgOneMultiWrong}{0.04}
\newcommand{\rgOneMultiArgmax}{0.84}
\newcommand{\rgOneStaleN}{2376}
\newcommand{\rgOneStaleLww}{0.92}
\newcommand{\rgOneStalePal}{0.51}
\newcommand{\rgOneStaleAbst}{0.45}
\newcommand{\rgOneStaleWrong}{0.04}
\newcommand{\rgOneStaleArgmax}{0.81}
\newcommand{\rgOneDelayN}{159}
\newcommand{\rgOneDelayLww}{0.08}
\newcommand{\rgOneDelayPal}{0.14}
\newcommand{\rgOneDelayAbst}{0.65}
\newcommand{\rgOneDelayWrong}{0.21}
\newcommand{\rgOneDelayArgmax}{0.49}
\newcommand{\rgOneLowN}{364}
\newcommand{\rgOneLowLww}{0.04}
\newcommand{\rgOneLowPal}{0.02}
\newcommand{\rgOneLowAbst}{0.92}
\newcommand{\rgOneLowWrong}{0.07}
\newcommand{\rgOneLowArgmax}{0.61}
\newcommand{\rgOneOthN}{234}
\newcommand{\rgOneOthLww}{0.06}
\newcommand{\rgOneOthPal}{0.03}
\newcommand{\rgOneOthAbst}{0.91}
\newcommand{\rgOneOthWrong}{0.06}
\newcommand{\rgOneOthArgmax}{0.37}
\newcommand{\rgOneRetrN}{463}
\newcommand{\rgOneRetrLww}{0.49}
\newcommand{\rgOneRetrPal}{0.35}
\newcommand{\rgOneRetrAbst}{0.35}
\newcommand{\rgOneRetrWrong}{0.30}
\newcommand{\rgOneRetrArgmax}{0.49}
\newcommand{\rgOneDisputedN}{2279}
\newcommand{\rgOneArgmaxFull}{0.593}
\newcommand{\rgOneLwwFull}{0.614}
\newcommand{\rgOneHalfCov}{0.49}
\newcommand{\rgOneHalfArgmax}{0.690}
\newcommand{\rgOneHalfLww}{0.638}
\newcommand{\rgPoisN}{78}
\newcommand{\rgPoisLww}{0.05}
\newcommand{\rgPoisPalWrong}{0.76}
\newcommand{\rgPoisArgmax}{0.20}
\newcommand{\rgPoisPalInj}{0.78}
\newcommand{\rgPoisLwwInj}{1.00}
\newcommand{\rgPoisPalAbst}{0.22}
\newcommand{\rgPoisArgmaxInj}{0.82}
\newcommand{\rgPoisPalWorld}{0.03}
\newcommand{\suFaSelfAbst}{0.321}
\newcommand{\suFaSelfWrong}{0.000}
\newcommand{\suFaSelfWorld}{0.679}
\newcommand{\suFaDelayAbst}{0.627}
\newcommand{\suFaDelayWrong}{0.267}
\newcommand{\suFaDelayWorld}{0.107}
\newcommand{\suFaLowAbst}{0.948}
\newcommand{\suFaLowWrong}{0.042}
\newcommand{\suFaLowWorld}{0.010}
\newcommand{\suFaOthAbst}{0.906}
\newcommand{\suFaOthWrong}{0.083}
\newcommand{\suFaOthWorld}{0.010}
\newcommand{\suFaSelfN}{53}
\newcommand{\suFaSelfLww}{0.962}
\newcommand{\suFaProp}{20.4}
\newcommand{\suFaPropLo}{18.0}
\newcommand{\suFaPropHi}{22.7}
\newcommand{\suFaWarrPerStream}{7.70}
\newcommand{\suFbSelfAbst}{0.000}
\newcommand{\suFbSelfWrong}{0.038}
\newcommand{\suFbSelfWorld}{0.962}
\newcommand{\suFbDelayAbst}{0.587}
\newcommand{\suFbDelayWrong}{0.280}
\newcommand{\suFbDelayWorld}{0.133}
\newcommand{\suFbLowAbst}{0.938}
\newcommand{\suFbLowWrong}{0.052}
\newcommand{\suFbLowWorld}{0.010}
\newcommand{\suFbOthAbst}{0.896}
\newcommand{\suFbOthWrong}{0.083}
\newcommand{\suFbOthWorld}{0.021}
\newcommand{\suFbSelfN}{53}
\newcommand{\suFbSelfLww}{0.962}
\newcommand{\suFbProp}{21.0}
\newcommand{\suFbPropLo}{18.4}
\newcommand{\suFbPropHi}{23.6}
\newcommand{\suFbWarrPerStream}{7.66}
\newcommand{\trPalCW}{0.063}
\newcommand{\trPalWC}{0.012}
\newcommand{\trPalWW}{0.075}
\newcommand{\trPalInst}{0.092}
\newcommand{\trLogCW}{0.080}
\newcommand{\trLogWC}{0.032}
\newcommand{\trLogWW}{0.681}
\newcommand{\trLogInst}{0.404}
\newcommand{\trTdbCW}{0.003}
\newcommand{\trTdbWC}{0.003}
\newcommand{\trTdbWW}{0.576}
\newcommand{\trTdbInst}{0.150}
\newcommand{\trN}{587}
\newcommand{\primN}{741}
\newcommand{\primEmpty}{21}
\newcommand{\primLogExact}{0.525}
\newcommand{\primLogExactNe}{0.540}
\newcommand{\ovStreams}{303}
\newcommand{\ovQueries}{21{,}825}
\newcommand{\ovDerived}{10{,}847}
\newcommand{\ovDisagree}{0}
\newcommand{\ovProductFail}{0}
\newcommand{\ovClosureFail}{0}
\newcommand{\ovMaxEnum}{15}
\newcommand{\ovSkipped}{1299}
\newcommand{\holmHThree}{0.998}
\newcommand{\holmHTwoR}{0.228}
\newcommand{\holmHTwoLo}{0.224}
\newcommand{\holmHTwoHi}{0.233}
\newcommand{\holmHThreeEst}{0.062}
\newcommand{\holmHThreeLo}{0.015}
\newcommand{\holmHThreeHi}{0.078}
\newcommand{\holmHFourA}{-0.492}
\newcommand{\holmHFourALo}{-0.533}
\newcommand{\holmHFourAHi}{-0.451}
\newcommand{\holmHFourB}{0.310}
\newcommand{\holmHFourBLo}{0.286}
\newcommand{\holmHFourBHi}{0.333}
\newcommand{\fieldProvTwo}{2.3}
\newcommand{\fieldProvThree}{5.3}
\newcommand{\spendTotal}{88.85}


\title{\pal: Accountable Belief Revision for LLM Systems\\
\large When does write-time justification maintenance beat query-time adjudication?}
\author{Christian Leiva Beltr\'an\thanks{Research design, semantics and experiments were developed with AI assistance (Claude); the author reviewed and takes responsibility for the content. Code, data generators, prompts, cost ledgers and the pre-registration are released with the paper.}}
\date{3 October 2026 --- preprint, v0.3 draft}

\begin{document}
\maketitle

\begin{abstract}
LLM systems that persist knowledge must revise it when observations conflict, arrive late, come from unreliable sources, or are withdrawn. We study one capability such systems lack: propagating the \emph{retraction} of evidence through the justifications of the conclusions that depended on it, while refusing to commit to a revision the evidence does not determine. We give a declared semantics of \emph{justified belief} for streams of observations (an assumption-based truth-maintenance oracle over a small claim language with explicit operator cues, dispute and withdrawal), a benchmark generator, \rs, that separates the world, the observations, and the beliefs those observations justify, and a representation-neutral output contract under which every system---including an append-only evidence log with query-time adjudication---is scored for preserving or prematurely resolving uncertainty. \pal{} is a write-time belief store whose derived beliefs pin the versions of the evidence that justified them; it is compared with query-time adjudication and with ablations at matched cost. In the structured setting (500 streams, no LLM) the store reproduces the oracle exactly, removing retraction propagation costs \sOnePropEffect{} points [\sOnePropLo, \sOnePropHi] on warranted downstream queries and nothing on direct ones, and a forced-commitment variant raises unsupported propagation by \sOneCommitUnsup{} points [\sOneCommitUnsupLo, \sOneCommitUnsupHi] without losing resolvable accuracy; a learned-prior evidence model is better calibrated than class priors alone (ECE \RqLearnedEce{} vs.\ \RqPriorEce) but worse than a logistic regression over the same evidence features (\RqLogregEce); an explicit attacker model shows that cue resolution creates an attack surface that provenance alone does not close. In the natural-language setting (100 streams, claims extracted by Claude Haiku 4.5 with a 1.5\% slot-error rate), \pal{} over the extracted claims scores \sTwoPalWarr{} on warranted downstream queries against \sTwoLogWarr{} for the same claims adjudicated by the LLM at query time, a paired difference of \sTwoHOne{} points [\sTwoHOneLo, \sTwoHOneHi], while committing on \sTwoPalCommit{} of unresolved slots against \sTwoLogCommit, at a third of the LLM cost; the sign holds with Claude Sonnet 4.6 (+\sStrongHOne{} points). Two pre-registered hypotheses fail: the evidence model's support is less well calibrated than the adjudicator's own confidence on extracted claims, and a controlled corruption curve shows justified-belief maintenance degrading faster than last-write-wins as extraction errors grow, because every extraction error becomes a standing dispute. On 30 streams assembled from Wikidata revision histories (pseudonymised, with the editors' deprecations as withdrawals) the propagation effect transfers ($+\sThreeProp$ points) and \pal{} matches the rule-based ceiling; the LLM adjudicator's gap there is decomposed into two output conventions and a conflict between its priors and the declared schema. On LongMemEval's knowledge updates, which score the latest self-report as truth, the semantics' refusal to assume recency costs it most of its coverage, which we report as a mismatch of scoring targets. We report the pre-registered hypotheses, every negative result, and measured costs (\$\spendTotal{} in total).
\end{abstract}

\section{Introduction}
Memory for LLM agents is usually built as retrieval over an append-only store, with an LLM adjudicating conflicts at query time. This works until the evidence itself changes: a source is withdrawn, an earlier report is corrected, a change is reported late. On the STALE benchmark most memory frameworks score below 10\% on implicit memory conflicts and the best standalone model reaches 55.2\% \citep{stale2026}; knowledge editors update a target fact but not its consequences \citep{rippleedits}; and finetuning on new facts erodes prior knowledge. The question we ask is narrower than ``how should agents learn continually'' and, we argue, prior to it: \emph{when evidence is withdrawn or corrected, which conclusions should lose support, and does maintaining that structure at write time earn its cost over re-deriving everything at query time?}

Two design commitments shape the study. First, systems are scored against \emph{justified} belief, not against the world: an agent that believes a true thing for no good reason is not rewarded, and an agent that holds two readings open when the evidence does not decide between them receives full credit. This requires a reference reasoner that sees only the observations and the declared assumptions and never the world, and that retains alternative interpretations rather than selecting one (\S\ref{sec:semantics}). Second, the requirement to preserve uncertainty is placed on observable behaviour through a shared output contract (\S\ref{sec:contract}), so that an append-only log adjudicated by an LLM competes on equal terms with a store that represents hypotheses explicitly.

Our contributions are: (1) a declared, total semantics of justified belief for observation streams with operator cues, dispute, correction and withdrawal, implemented by two independent oracles that agree on every generated stream; (2) \rs, a generator with a perturbation menu (delays, unknown change dates, reversals, correlated copies, errors from trusted sources, genuine low-trust corrections, self-corrections, explicit retraction of reports and of sources, an attacker) and a declared stratified query selection; (3) \pal, a write-time store with versioned keys and pinned derived justifications, and the measurement that retraction propagation through justifications---not ordinary belief updating---is what prevents downstream errors; (4) an evidence model with a defined meaning of support, evaluated for calibration on held-out streams against learned baselines, with an honest negative result; (5) the Setting 2 result: under automatically extracted dependencies, write-time justification maintenance beats LLM query-time adjudication on warranted retraction consequences by a wide margin at lower cost, and the regime in which it stops paying (extraction error rates of tens of percent) is measured. Retraction propagation as a \emph{mechanism} is not new: MemTX \citep{memtx2026} cascades typed repair through derived records in transactional multi-agent memory, and CUPMem \citep{stale2026} invalidates dependent attributes at write time; our claim is about when it is warranted, what it costs, and how to measure it without rewarding premature commitment.

\section{Related work}
\paragraph{Belief revision and truth maintenance.} AGM revision \citep{agm1985} and belief-base postulates \citep{hansson1999} define rational change of a theory; Katsuno and Mendelzon \citep{km1991} distinguish \emph{update} (the world changed) from \emph{revision} (we were wrong), a distinction our operator cues make observable. Justification- and assumption-based truth maintenance \citep{doyle1979,dekleer1986} maintain which conclusions remain supported when premises are withdrawn; our oracle is an ATMS over environments of accepted reports, and \pal{} is a JTMS-like store with versioning. Non-prioritised revision \citep{hansson1997,makinson1997} allows input to be rejected; our admission policy is non-prioritised in that reliability never resolves disputes at the oracle level.
\paragraph{Temporal and bitemporal memory.} Bitemporal databases \citep{snodgrass} separate valid time from transaction (belief) time; Kumiho \citep{kumiho2026} and TOKI \citep{toki2026} bring versioned revisions and bitemporal operators to agent memory, Zep/Graphiti \citep{graphiti2025} build temporal knowledge graphs. None evaluates retraction propagation under extracted dependencies against query-time adjudication.
\paragraph{Benchmarks.} STALE \citep{stale2026}, LongMemEval \citep{longmemeval}, LoCoMo and StateMemBench test memory updates and state tracking against ground truth; \rs{} differs in scoring against justified belief with hidden world state, delayed and unreliable observations, and evidence withdrawal.
\paragraph{Calibration and source reliability.} Nous \citep{nous2026a,nous2026b} studies reliability-conditioned updating and shows that learning useful state decisions can need quadratically fewer records than calibrating sources; we take from it the choice of decision-relevant metrics (risk--coverage) alongside Brier and ECE, and class-level rather than per-source priors.

\section{Semantics of justified belief}\label{sec:semantics}
We fix a small claim language and define, for any stream of observations, the set of \emph{admissible interpretations} and the beliefs they justify. Everything in this section is a declared benchmark assumption; the oracle never sees the world.

\paragraph{Claim language.} A domain declares entities $E$, attributes, and rules. Each attribute $a$ declares a value type, a cardinality (\texttt{single} or \texttt{multi}), whether it is \texttt{changeable} (volatile vs.\ stable), whether a report of it may be hypothesised erroneous (\texttt{error\_allowed}), and whether differing values on one key are evidence against each other (\texttt{competing\_values}; two residences compete, two nicknames coexist). A proposition is $(e,a,v)$ on key $(e,a)$. Rules are Horn clauses, \texttt{strict} or \texttt{defeasible} with declared exceptions; heads are derived attributes (e.g.\ $\mathit{work\_city}(e)\leftarrow \mathit{employer}(e,x)\wedge \mathit{hq\_city}(x,c)$ unless $\mathit{remote}(e)$).

\paragraph{Time and observations.} Valid time is an integer day; belief time is report time. An \emph{assertion} is $(\mathit{id}, t_{\mathrm{rep}}, \mathit{source}, e,a,v, \mathit{valid\_cue}, \mathit{op\_cue})$ where the valid cue is \texttt{none}, \texttt{since}~$t$ or \texttt{at}~$t$ and the operator cue is \texttt{none}, \texttt{change} (optionally \emph{from} $v_{\mathrm{prev}}$) or \texttt{correction} (\emph{of} a report id). A \emph{retraction} withdraws a report or every report of a source for all later belief times. Sources declare a reliability class (\texttt{trusted}, \texttt{standard}, \texttt{low}, \texttt{blocked}) and an \emph{origin}; sources sharing an origin are copies. A report with no temporal cue is taken to assert that the value held at least at reporting time (A1); \texttt{since}~$t$ anchors it at $t$ and \texttt{at}~$t$ asserts it at $t$. Absent later evidence of a different value, a value persists from its anchor (A4, inertia); the valid cue fixes when a value is known to hold, never when it stops, so inertia applies to \texttt{at} reports exactly as to the others.

\paragraph{Admission and withdrawal.} Policy P0: every report from a non-blocked source is admitted; reliability classes never affect admissibility, only the support layer. A correction whose source shares the origin of the corrected report \emph{withdraws} that report (A-SELF): an origin's later statement about its own earlier statement is a withdrawal, not a competing report. A correction from another origin is a competing report with a cue.

\paragraph{Dispute and contradiction.} Two admitted assertions on the same key are \emph{disputed} iff they have different values and the attribute declares competing values. A \emph{contradiction} between two accepted assertions arises only for single-valued keys when the attribute is stable, or when the attribute is changeable and the two are anchored at the same time. Dispute licenses the hypothesis of error; contradiction forbids joint acceptance. The distinction is what makes ``the London report may be wrong'' admissible even when London-then-Paris is consistent as a change.

\paragraph{Interpretations.} At belief time $\tau$, an interpretation labels every admitted assertion \textsc{true} or \textsc{err}, and every maximal run of accepted same-valued reports on a multi-valued changeable key that is followed by a different value \textsc{cont} or \textsc{ended}. It is admissible iff:
(A-ERR) a report may be \textsc{err} only if error is allowed and some \emph{accepted} competitor disputes it, or an accepted correction targets it---every error label is explained by an accepted competing report, and ``everything is wrong'' is never admissible;
(A-CORR) an accepted correction forces its target \textsc{err};
(A-CHG) an accepted change \emph{from} $v_{\mathrm{prev}}$ presupposes the previous value: if any report of $v_{\mathrm{prev}}$ is anchored strictly earlier, at least one must be accepted (a presupposition about the value, not about a particular report), and on a multi-valued key the run of $v_{\mathrm{prev}}$ preceding the change is labelled \textsc{ended} (the \emph{from} clause states that the previous value stopped);
(A-CUE, policy P0c) a report carrying a change cue is not disputed by competitors anchored strictly earlier than itself---the cue identifies the revision operator, so the earlier value is explained rather than contested---while same-time and later competitors still dispute it;
(A-CONS) no contradiction among accepted reports.
Minimality is deliberately \emph{not} imposed: preferring fewer error labels would make ``London$\to$Paris was a move'' established and ``the London report was wrong'' inadmissible, which is exactly the premature commitment the benchmark is designed to penalise. Preferences belong to the support layer. An interpretation is identified by its \textsc{err} set and the timelines it induces; keys are independent, so interpretation sets are computed per key, and the oracle is made total by a declared relaxation ladder that is counted whenever it fires (it never fires on generated streams under P0c).

\paragraph{Timelines and answers.} Under an interpretation, each accepted run holds from its first anchor; on single-valued keys it ends at an unknown change point in the gap before the next different accepted anchor; on multi-valued keys it persists (\textsc{cont}) or ends in that gap (\textsc{ended}). The value at $t$ inside a gap is $\{v_{\mathrm{prev}}\}$ or $\{v_{\mathrm{next}}\}$---unknown change dates are uncertainty by construction. Derived attributes are forward-chained per interpretation, branching over candidate sets; a defeasible rule whose exception holds contributes nothing.

\paragraph{Gold.} For a value slot, the candidates are the union over interpretations: one candidate $\Rightarrow$ \texttt{established}; several $\Rightarrow$ \texttt{unresolved} with exactly that set; none $\Rightarrow$ \texttt{unknown}. For a multi-valued slot a candidate is a \emph{set} of values and alternatives are sets (London-only, Paris-only and both are three alternatives, never flattened), and the empty set is a candidate like any other: when every interpretation justifies no member---nothing admitted on the key, or everything withdrawn---the slot is \texttt{established} with $\emptyset$, which reads ``no affiliation is justified'' and is distinct from \texttt{unknown} (no interpretation at all, or an undeclared key). A yes/no proposition true in every interpretation is \texttt{established} (false in every one: established with negative polarity---the semantics is closed over the admitted evidence, so a value never asserted does not hold); mixed $\Rightarrow$ \texttt{possible}. These three conventions (empty set, established-false, unknown) are part of the contract and, as \S\ref{sec:res-s3} reports, the one for empty sets was missing from the instruction the LLM adjudicators received. A downstream query asked right after a retraction is a \emph{warranted-retraction consequence} iff its gold is established and differs from its gold one belief step earlier; only such queries enter the primary endpoint.

\paragraph{Worked case.} With \emph{residence} declared multi-valued and competing, after ``Alex lives in London'' (o1) and ``Alex lives in Paris'' (o2) with no cues, exactly four interpretations are admissible: both accepted with London continuing (both residences), both accepted with London ended (a move), o1 in error, o2 in error. ``Did Alex move?'' is \texttt{possible}; ``Was the London report wrong?'' is \texttt{possible}; ``Which residences were reported?'' is \texttt{established} with provenance; ``Which city's tax applies?'' (a rule on residence) is \texttt{unresolved} over three alternatives and is not a warranted consequence. Had o2 said ``Alex moved from London to Paris'', the move would be established under A-CUE. Two independent implementations of this semantics (a constructive per-key enumeration and a brute-force global one) were compared on the three hand-written reference streams, where they agree on the interpretation sets and the gold of every query at every belief time; on the 500 generated streams only the weaker check that every key has at least one admissible interpretation under P0c was run (the global oracle does not scale to them). The incremental store and the symbolic replay, which share the per-key oracle, are compared on every query of every setting in \S\ref{sec:res-systems}.

\paragraph{Self-update (A-SU, variant P0cSU).} A later uncued report from the \emph{same origin} on a single-valued changeable key is, under P0c, a dispute with the origin's own earlier report. The regime analysis of \S\ref{sec:res-regime} and an external check showed this to be a poor default for self-report streams. The Bayesian argument depends on the observation model, and we state both that we considered. Let two reports of one source on one key differ in value; let $c$ be the probability that the world value changed between them and $e$ the probability that a report is erroneous, independently. Under a \emph{binary} value model (an erroneous report shows the only other value), $P(\text{later value is current}\mid\text{values differ})=\frac{c(1-e)^2+(1-c)e(1-e)}{c[(1-e)^2+e^2]+2(1-c)e(1-e)}$, which exceeds $\tfrac12$ by $c(1-2e)/(2\cdot\text{denominator})$, i.e.\ for every $c>0$ and $e<\tfrac12$ (an earlier draft omitted the $ce^2$ term, which inflated the value at large $c$). Under an \emph{open} value model (an erroneous report shows a value distinct from both the true ones), the same probability is $\frac{c(1-e)+(1-c)e(1-e)}{c+(1-c)(2e-e^2)}$, which exceeds $\tfrac12$ only when $c(1-2e)\ge(1-c)e(1-e)$, roughly when changes are at least as likely as errors. In either model the later value is the better \emph{guess} in the stated regime; a probability above one half supports a preference, not a justification, so A-SU is a declared policy choice and not a theorem. A-SU, pre-registered before it was implemented (Appendix~\ref{app:chronology}), forbids labelling a report \textsc{err} when every accepted competitor is a strictly-earlier report from its own origin and no correction targets it; it never defeats a later competitor, a correction or a withdrawal, and does not apply to stable attributes. It is adopted as the default for self-report streams and reported as the declared variant P0cSU elsewhere; its measured cost is in \S\ref{sec:res-regime}; the frozen results of Settings 1--3 are under P0c.

\paragraph{Variants.} Literal P0 (no cue shields any report: only withdrawal resolves) and P1 (class-prioritised dispute: a lower-class report disputes a higher-class one only when corroborated from a second origin) are implemented and reported as sensitivity analyses. P1 is not the main configuration because it builds a reliability ordering into the oracle that the systems under test are meant to estimate, and because it marks a trusted source's false statement \texttt{established} whenever only lower-ranked sources contradict it.


\section{The output contract}\label{sec:contract}
The requirement to preserve uncertainty is placed on \emph{observable behaviour}, not on internal representation: the scoring must not favour a system that stores hypotheses over one that reasons at query time. Every system, including the references, answers through one JSON object: a \texttt{status} in \{\texttt{established}, \texttt{possible}, \texttt{unresolved}, \texttt{unknown}\}, an \texttt{assertion}, an \texttt{alternatives} list, and \texttt{provenance} (observation ids); an optional free-text explanation is never scored. The statuses are defined operationally and are mutually exclusive: \texttt{established} means supported under every admissible interpretation given the declared assumptions, never independently verified truth; \texttt{possible} qualifies one proposition that holds under some interpretations; \texttt{unresolved} declines to choose between candidate values and must list them; \texttt{unknown} has no candidates. The shared instruction contains these definitions and the Alex examples verbatim and is identical for every system, as is the vocabulary (entity ids, attribute types and normalisation, observation ids). Nothing exposes gold labels or hidden dependency rules beyond what the setting declares.

\paragraph{Deterministic scoring.} Values are canonical (entity ids, ISO dates, lower-cased strings; sorted lists for multi-valued slots), so credit is exact: status must match, and the assertion or the alternative set must equal the gold after normalisation. Selecting one interpretation on an \texttt{unresolved} or \texttt{possible} slot scores zero on that query and counts as a \emph{premature commitment}. A lenient variant (superset of alternatives) was defined but is not reported in this version; all numbers are exact-match. No LLM judge is used in Settings 1--2.

\paragraph{Provenance.} A cited observation id is \emph{supporting} when it asserts a (key, value) binding that lies on an oracle justification path for the answer: for a derived slot the rule engine records the bindings used along every derivation that yields a candidate value. Provenance precision and recall are secondary metrics; an existing but non-supporting id earns nothing.

\paragraph{Malformed outputs.} Outputs are validated against the schema and vocabulary; exactly one repair attempt with a fixed, system-independent prompt is allowed and charged to the system; a second failure is scored wrong on every query (never as abstention) and the malformed rate is reported (it is zero for every system in every run; the repair was invoked once in the Haiku runs and on 192 of 994 answers in the Sonnet subset).

\paragraph{References.} Always-abstain answers every value slot \texttt{unresolved} over the reported values (and yes/no slots \texttt{possible}); always-commit takes the most recent report at face value. Neither is a competitor: a real system must beat always-abstain on resolvable queries and always-commit on unresolved ones, otherwise uncertainty handling has become an escape from answering, or commitment an escape from reasoning.


\section{\rs}\label{sec:benchmark}
\rs{} separates three layers per stream: the world trajectory (hidden from every system, used only for calibration and for the attacker), the observation stream with provenance (what systems see), and the oracle's justified beliefs (gold, computed from the observations alone).

\paragraph{World and reporting.} A seeded simulator draws persons, organisations and cities with attributes \emph{employer, residence (multi-valued), remote, birth\_date, nickname (non-competing), hq\_city, ceo}, three rules (\emph{work\_city, local\_tax\_city, ceo\_city}; one chained, one with a defeasible exception), and changes of volatile attributes over a 60-day horizon with occasional reversals. Eight sources in four classes report each world fact one to three times with geometric delays and a heavy tail (10--25 days), with an explicit \texttt{since} cue half the time on changeable attributes, with change cues on a third of change reports (half of them with a stated previous value), and with class-specific error rates (trusted 6\%, standard 12\%, low 30\%). Two sources are copies of others (same origin). An erroneous report is followed with declared probabilities by a self-correction (A-SELF), by a cross-origin correction from a random source that is low-trust half the time (a genuine low-trust correction), or by nothing. Reports are explicitly retracted at a rate that favours erroneous ones, and one non-trusted source is retracted wholesale in a third of streams. Per-key report counts are capped at seven to bound enumeration. The attacker model (secondary RQ) injects false assertions at a declared rate through a fresh low-trust channel or a compromised trusted channel, half of them with a false change cue.

\paragraph{Queries.} Event-driven queries are asked at every retraction, correction and cued change: the affected key (\texttt{current}), every derived key whose justification can depend on it (\texttt{downstream}, found by inverting the oracle's own dependency relation), the relevant yes/no propositions (\emph{erroneous}, \emph{changed}), each with a companion asked one belief step earlier. Grid queries sample the bitemporal cells: as-of valid time, as-of belief time (``what was justified on day $\tau'$''), \emph{holds}, and \texttt{reported}. The declared selection keeps at most 44 primary queries per stream (a cap, not a guaranteed count: a stratum shorter than its quota is not padded): 35\% warranted-retraction consequences (downstream first), 25\% non-warranted event queries (they feed the instability endpoint), 40\% grid queries balanced between resolvable and unresolved, plus the $\tau-1$ companions of event queries and the \texttt{reported} slots, which are scored separately. Generated streams average 131 observations and about 280 candidate queries before selection; the frozen Setting 1 set has 60.5 queries per stream (37.1 primary, 20.5 companions, 2.9 \texttt{reported}; 30{,}272 in all), of which \sOneWarrPerStream{} per stream are warranted downstream queries; the frozen Setting 2 set has 67.4 per stream (40.6 primary, 23.4 companions, 3.4 \texttt{reported}; 6{,}737 in all) and \sTwoWarrPerStream{} warranted downstream queries per stream (741 in total).

\paragraph{Natural-language setting.} Each observation is rendered as a report line ``[id] (outlet, date): sentence'' by templates with fictional names (so no world knowledge can be leaned on), paraphrased by a cheap model, and passed through a cue-recovery check---deterministic markers plus a cheap-model pass blind to the structured source---with up to two regenerations; lines that still fail are adjudicated by stripping the lost cue from the \emph{frozen label}, so gold reflects the text as presented. The evaluation set is then frozen. A shared extractor turns every line into a structured claim under the declared attribute and rule schema (the rules are given, not extracted); its slot errors, \emph{missed ambiguity} (text has no cue, extractor added one) and \emph{spurious ambiguity} (text has a cue, extractor dropped it) are measured against the frozen labels, and an extractor's declaration never changes which queries are answerable.


\section{Systems}\label{sec:systems}
\paragraph{\pal.} A write-time belief store. Every observed key holds an append-only list of versions (belief time, admitted report ids, interpretation set computed with the declared semantics); every derived belief (entity, derived attribute) holds versions that \emph{pin} the key versions it was justified by. Answers at belief time $\tau$ are read from the versions in force at $\tau$, so bitemporal questions are lookups. On an assertion, the key is recomputed and every dependent derived belief is refreshed (ordinary belief updating). On a retraction or withdrawal, the key is recomputed and---this is the propagation step---every derived belief that depends on it is refreshed. Support (\S\ref{sec:support}) is computed from the admitted evidence of a key. In the natural-language setting the store reasons over the shared extracted claims; the LLM is used only at the boundary (extraction), never to answer.

\paragraph{Ablations.} \emph{No propagation}: a retraction updates the retracted key only; derived beliefs keep their pinned versions until an ordinary assertion refreshes them. \emph{Forced commit}: each key keeps a single interpretation, the admissible one with fewest error labels (ties: ended over continuing). Further ablations of the support model are reported with RQ3.

\paragraph{Symbolic replay (the equivalence baseline).} The same claims in an append-only log; at query time the semantics is executed from scratch over the log filtered to the query's belief time: every key the query depends on is re-enumerated and the derived value re-chained, with the per-key cache cleared before each query (\emph{cold}) or retained across queries (\emph{cached}): the cache is a memo keyed by (key, belief time, policy) that is never invalidated, which is sound because the log is append-only and a report at a later belief time cannot change an earlier belief time's interpretation set; it therefore helps only when a later query touches the same key at the same belief time. Replay and \pal{} execute the same declared semantics on the same evidence and have no information the other lacks; they must agree on every answer, and \S\ref{sec:res-systems} verifies that they do and measures what each costs. In the natural-language settings replay is the ``log + rule adjudicator'' row: what query-time reasoning achieves over the extracted claims when the reasoning is correct.

\paragraph{Append-only log with LLM adjudication.} The same extracted claims in the same log; at query time the backbone LLM receives the log filtered to the query's belief time, the declared attribute and rule block, the shared instruction, and the query rendered as its answer slot, and must return the contract JSON, with one repair attempt. Context-budget variants restrict the log to the $k$ most recent entries plus every entry on the queried key ($k\in\{10,20,40\}$); a raw-report variant adjudicates the natural-language reports directly, with no extraction. Relative to \pal{} this comparison changes two things at once, who reasons (a prompted model instead of a rule engine) and when (query time instead of write time); replay, which reasons with the rule engine at query time, separates them.

\paragraph{Temporal DB (last write wins).} A rule-based bitemporal store without justifications: the record with the latest valid-time anchor wins, retractions remove records, no error hypotheses, derived values computed at query time. It differs from \pal{} in two respects at once, no justifications and a commit-always uncertainty policy, so its gap is not attributed to one of them; the always-commit reference (latest report by report time) isolates pure recency.

\paragraph{Support model.}\label{sec:support} Support is the probability that a candidate value is true for its interval, defined from admitted evidence only: evidence is grouped by origin (copies count once), each origin contributes a per-class log-odds (declared, or learned from verified outcomes of training streams), discounted by staleness on changeable attributes with a 20-day half-life, boosted by a change cue, and capped for a report targeted by a correction at half the correcting origin's weight (the min(provenance, content) cap); support of $v$ is the logistic of the net evidence for $v$ against all competitors, with a prior offset when no competitor exists. Source identities are never used, only classes and origins. Support is produced for exactly the alternatives the justified-belief layer keeps open, so the two layers are evaluated separately.


\section{Experimental protocol and pre-registration}\label{sec:protocol}
\paragraph{Endpoints.} The primary endpoint is exact-match accuracy on warranted-retraction consequence queries on \emph{downstream} slots (propagation through justifications); the same on the retracted key itself is secondary (ordinary updating). Ambiguity is measured separately and never pooled with the primary: the premature-commitment rate on unresolved/possible slots, and the answer-instability rate (registered as ``unsupported propagation'': share of non-warranted event queries whose answer changed between the companion and the event while the gold did not), decomposed post hoc into correct/wrong transitions because the registered form counts beneficial changes as failures, reported next to resolvable accuracy and the two references so that abstention cannot buy a low commitment rate. Calibration is claim-level Brier, ECE (15 equal-mass bins) and the area under the risk--coverage curve. Cost is input plus $8\times$ output tokens (the backbone's price ratio) summed over ingestion, repair and answering, with contract tokens charged at full price. The stream is the unit of analysis: paired per-stream differences with a 10{,}000-resample cluster bootstrap and effect sizes with 95\% intervals.

\paragraph{Pre-registered hypotheses.} H1 (primary): in Setting 2 on warranted downstream queries, \pal{} exceeds the LLM-adjudicated log at matched cost by $\geq 10$ points with an interval excluding zero; falsified if the interval includes zero under extracted dependencies, unfamiliar sources and matched cost together, in which case the cost-crossover analysis is the main finding. H2: a crossover query-to-update ratio $r^*<1$ exists. H3: the evidence model's support has lower ECE than verbalised confidence on held-out sources. H4 (joint): lower premature-commitment rate than the log baseline \emph{and} resolvable accuracy not lower by more than 3 points. H2--H4 use Holm correction (\S\ref{sec:threats} reports why it changes no conclusion). H1's ``matched cost'' clause could not be executed as registered; \S\ref{sec:res-s2} reports the context-budget curve that replaces it. Appendix~\ref{app:chronology} gives the dated chronology of every decision, including those taken after results were seen. Phase gates preceded any paid run: two independent oracles must agree; a propagation-off store must score measurably worse on warranted downstream queries; the references must split; a forced-commit store must raise unsupported propagation. The pre-registration, with every deviation dated, is released with the code.

\paragraph{Deviations recorded.} (i) The planned backbone (GPT-5 mini via a batch API) and Amazon Bedrock were unreachable from the execution environment; Claude Haiku 4.5 through the Anthropic Message Batches API (50\% discount, prompt caching on the shared instruction) was used for every stage, and Claude Sonnet 4.6 on 15 streams as the sign-stability check. (i$'$) Setting 2 ran on 100 streams rather than 150 to keep total spend under the \$100 ceiling (interval half-width 0.04 at the observed spread); the raw-report log and the matched-cost variant ran on 30 streams. Total measured spend for every LLM call in the study, including development, two discarded batches and the post-review reruns, was \$\spendTotal. (ii) The warranted-query floor of 8 per stream was met on development streams by raising the primary query budget from 40 to 44 with quota 35/25/40 before the freeze; the frozen Setting 1 set has \sOneWarrPerStream{} and the frozen Setting 2 set \sTwoWarrPerStream{} (741 queries; 99 of 100 streams contribute to the primary estimate); generator rates were not changed. (iii) Query counts (\S\ref{sec:benchmark} gives the per-stream flow): the \sOneAnswers{} ``scored answers'' of Setting 1 are the primary queries whose gold is \texttt{established}, \texttt{unresolved} or \texttt{possible} (companions, \texttt{reported} slots and \texttt{unknown}-gold queries excluded); the Setting 2 totals are 6{,}737 queries, 4{,}062 primary, 741 warranted downstream; the literal-P0 yield of 2.3 warranted downstream queries per stream was measured on development streams before the selection rule, the 3.38 of the sensitivity table after it. (iv) The dispute semantics changed from literal P0 to P0c after Phase 1 measurement (literal P0 leaves 2.3 warranted downstream queries per stream and 48\% unresolved slots), with literal P0 retained as a sensitivity variant; and the change presupposition (A-CHG) was weakened from a specific report to the value after it produced inconsistent streams.


\section{Results}\label{sec:results}
\subsection{Setting 1: structured streams, gold dependencies}
\begin{sidewaystable}\centering\footnotesize\setlength{\tabcolsep}{3.5pt}
\caption{Setting 1 (structured streams, gold dependencies, no LLM; 500 streams). Stream-level means with 95\% cluster-bootstrap intervals.}\label{tab:s1}
\begin{tabular}{lcccccccc}
\toprule
System & Warr.\ downstream & Warr.\ direct & Resolvable & Unresolved & Commit rate$\downarrow$ & Unsupp.\ prop.$\downarrow$ & Prov.\ P & Prov.\ R \\
\midrule
\textsc{Palimpsest} & 1.000 {\tiny[1.00,1.00]} & 1.000 {\tiny[1.00,1.00]} & 1.000 {\tiny[1.00,1.00]} & 1.000 {\tiny[1.00,1.00]} & 0.000 {\tiny[0.00,0.00]} & 0.000 {\tiny[0.00,0.00]} & 0.990 {\tiny[0.99,0.99]} & 0.993 {\tiny[0.99,0.99]} \\
\quad no propagation & 0.794 {\tiny[0.78,0.81]} & 1.000 {\tiny[1.00,1.00]} & 0.933 {\tiny[0.93,0.94]} & 0.962 {\tiny[0.96,0.97]} & 0.002 {\tiny[0.00,0.00]} & 0.017 {\tiny[0.01,0.02]} & 0.981 {\tiny[0.98,0.98]} & 0.989 {\tiny[0.99,0.99]} \\
\quad forced commit & 1.000 {\tiny[1.00,1.00]} & 1.000 {\tiny[1.00,1.00]} & 1.000 {\tiny[1.00,1.00]} & 0.042 {\tiny[0.04,0.05]} & 0.776 {\tiny[0.76,0.79]} & 0.090 {\tiny[0.08,0.10]} & 0.992 {\tiny[0.99,0.99]} & 0.775 {\tiny[0.77,0.78]} \\
Log + rule adjudicator & 1.000 {\tiny[1.00,1.00]} & 1.000 {\tiny[1.00,1.00]} & 1.000 {\tiny[1.00,1.00]} & 1.000 {\tiny[1.00,1.00]} & 0.000 {\tiny[0.00,0.00]} & 0.000 {\tiny[0.00,0.00]} & 1.000 {\tiny[1.00,1.00]} & 1.000 {\tiny[1.00,1.00]} \\
Temporal DB & 0.807 {\tiny[0.79,0.82]} & 0.909 {\tiny[0.89,0.92]} & 0.907 {\tiny[0.90,0.91]} & 0.000 {\tiny[0.00,0.00]} & 0.848 {\tiny[0.84,0.86]} & 0.105 {\tiny[0.09,0.12]} & 0.950 {\tiny[0.95,0.95]} & 0.341 {\tiny[0.33,0.35]} \\
Always abstain & 0.000 {\tiny[0.00,0.00]} & 0.000 {\tiny[0.00,0.00]} & 0.000 {\tiny[0.00,0.00]} & 0.393 {\tiny[0.38,0.41]} & 0.000 {\tiny[0.00,0.00]} & 0.075 {\tiny[0.06,0.09]} & 0.507 {\tiny[0.50,0.51]} & 0.739 {\tiny[0.73,0.75]} \\
Always commit & 0.438 {\tiny[0.42,0.46]} & 0.970 {\tiny[0.96,0.98]} & 0.770 {\tiny[0.76,0.78]} & 0.000 {\tiny[0.00,0.00]} & 0.793 {\tiny[0.78,0.81]} & 0.157 {\tiny[0.14,0.17]} & 0.939 {\tiny[0.93,0.94]} & 0.284 {\tiny[0.28,0.29]} \\
\bottomrule
\end{tabular}
\end{sidewaystable}
Table~\ref{tab:s1} reports all systems on 500 streams (\sOneAnswers{} scored answers per system). \pal{} reproduces the oracle on every query, so the incremental, versioned maintenance is correct and the bitemporal slots (as-of valid time, as-of belief time) are answered from stored versions without recomputation. Removing retraction propagation costs $\sOnePropEffect$ points $[\sOnePropLo, \sOnePropHi]$ on warranted downstream queries (paired over \sOnePropN{} streams) and nothing on direct ones---the signature the hypothesis predicts: ordinary belief updating handles the retracted key, only propagation handles its consequences. The forced-commit variant keeps perfect resolvable accuracy while committing on \sOneCommitRate\% of unresolved slots and changing downstream answers after non-warranted events $\sOneCommitUnsup$ points $[\sOneCommitUnsupLo, \sOneCommitUnsupHi]$ more often: premature commitment and unsupported propagation are measured as separate things and move independently of resolvable accuracy. The references split as required (always-abstain: \sOneAbstainRes{} resolvable / \sOneAbstainUnres{} unresolved; always-commit: \sOneCommitRes{} / \sOneCommitUnres{}). The temporal DB, which has valid time but no justifications, loses \sOneTdbWarr{} points on warranted downstream queries and commits on \sOneTdbCommit\% of unresolved slots. Provenance precision and recall of \pal{} are \sOneProvP/\sOneProvR{} against justification paths; the references cite existing but non-supporting ids. These are validation results on a benchmark whose generator and system share an author: they establish that the instruments discriminate, not that H1 holds.

\begin{table}[t]\centering\small
\caption{Setting 1: exact accuracy by query slot (mean over 500 streams).}\label{tab:slots}
\begin{tabular}{lcccccc}
\toprule
Slot & \textsc{Palimpsest} & no propagation & forced commit & Temporal DB & Always commit & Always abstain \\
\midrule
as-of valid time & 1.000 & 1.000 & 0.743 & 0.653 & 0.638 & 0.168 \\
as-of belief time & 1.000 & 1.000 & 0.808 & 0.777 & 0.782 & 0.223 \\
current & 1.000 & 1.000 & 0.723 & 0.675 & 0.701 & 0.242 \\
downstream (derived) & 1.000 & 0.870 & 0.664 & 0.578 & 0.296 & 0.000 \\
reported & 1.000 & 1.000 & 1.000 & 1.000 & 1.000 & 1.000 \\
changed? & 1.000 & 1.000 & 0.932 & 0.919 & 0.919 & 0.068 \\
erroneous? & 1.000 & 1.000 & 0.704 & 0.001 & 0.001 & 0.296 \\
holds at $t$? & 1.000 & 1.000 & 0.877 & 0.771 & 0.734 & 0.224 \\
\bottomrule
\end{tabular}
\end{table}
Table~\ref{tab:slots} breaks accuracy down by slot: the systems without justifications fail on \emph{erroneous?} almost completely (they never hold an error hypothesis), lose most on downstream slots, and answer as-of-belief-time questions from the current state rather than from what was justified then.

\begin{table}[t]\centering\small
\caption{Sensitivity of the Setting 1 gate to the dispute semantics (500 streams). Warranted downstream queries per stream, and the propagation on$-$off difference on them.}\label{tab:sens}
\begin{tabular}{lcccc}
\toprule
Policy & warr./stream & \textsc{Palimpsest} & no propagation & Temporal DB \\
\midrule
P0c (main) & 7.42 & 1.000 {\tiny[1.00,1.00]} & 0.794 {\tiny[0.78,0.81]} & 0.807 {\tiny[0.79,0.82]} \\
literal P0 & 3.38 & 1.000 {\tiny[1.00,1.00]} & 0.638 {\tiny[0.61,0.67]} & 0.918 {\tiny[0.90,0.93]} \\
P1 & 3.49 & 1.000 {\tiny[1.00,1.00]} & 0.688 {\tiny[0.66,0.72]} & 0.866 {\tiny[0.84,0.89]} \\
\bottomrule
\end{tabular}
\end{table}
Table~\ref{tab:sens} shows the sensitivity to the dispute semantics. Under literal P0 only withdrawal resolves, so a stream yields \sOneLitPZeroWarr{} warranted downstream queries and the propagation effect is larger ($+\sOneLitPZeroEffect$ points) because the surviving warranted queries are all retraction-driven; under P1 (class-prioritised dispute with corroboration) \sOnePOneWarr{} per stream and $+\sOnePOneEffect$ points. The qualitative conclusion does not depend on the policy; the size of the usable stratum does.

\subsection{RQ3: calibrated support}
\begin{table}[t]\centering\small
\caption{RQ3, Setting 1 (held-out streams): claim-level calibration of support against the hidden world. Lower is better.}\label{tab:rq3}
\begin{tabular}{lccc}
\toprule
Confidence & Brier & ECE & AURC \\
\midrule
Evidence model, declared priors & 0.182 & 0.118 & 0.138 \\
Evidence model, learned class priors & 0.175 & 0.113 & 0.131 \\
\quad without origin grouping & 0.179 & 0.121 & 0.136 \\
\quad without recency & 0.236 & 0.202 & 0.156 \\
\quad without cue adjustments & 0.183 & 0.128 & 0.136 \\
Class prior only & 0.255 & 0.227 & 0.168 \\
Logistic regression on evidence features & 0.138 & 0.034 & 0.112 \\
Constant 0.5 & 0.250 & 0.496 & 0.268 \\
\bottomrule
\end{tabular}
\end{table}
Support is evaluated at the claim level against the hidden world on 200 held-out streams (\RqN{} candidate values from unresolved and established value slots), with class priors learned on 300 training streams. The evidence model with learned priors is well calibrated (ECE \RqLearnedEce, Brier \RqLearnedBrier) and far better than class priors alone (\RqPriorEce, \RqPriorBrier); recency is the largest single contributor (without it ECE \RqNoRecEce), origin grouping the smallest (\RqNoOrigEce). It is, however, worse than a logistic regression over the same evidence features trained on the same streams (ECE \RqLogregEce, Brier \RqLogregBrier, AURC \RqLogregAurc{} vs.\ \RqLearnedAurc). We report this as a negative result for the hand-designed combination rule: the evidence model's value is its defined meaning and its auditability, not its calibration advantage. H3 proper, against verbalised and sampling-based confidence, is tested in Setting 2 below.

\subsection{Poisoning under an explicit attacker model}
\begin{table}[t]\centering\small
\caption{Poisoning (Setting 1, 60 streams per cell). Attacker injects false assertions at a rate of the honest ones, half with a false change cue; harm = share of value queries answered \emph{established} with the poisoned value; support = mean support the evidence model assigns to poisoned vs.\ honest candidates.}\label{tab:poison}
\begin{tabular}{llcccccc}
\toprule
Channel & Rate & \multicolumn{2}{c}{\textsc{Palimpsest}} & \multicolumn{2}{c}{Temporal DB} & \multicolumn{2}{c}{Support} \\
 & & harm$\downarrow$ & unresolved & harm$\downarrow$ & unresolved & poisoned & honest \\
\midrule
low & 0.01 & 0.003 & 0.27 & 0.003 & 0.03 & 0.21 & 0.67 \\
low & 0.05 & 0.024 & 0.28 & 0.030 & 0.03 & 0.33 & 0.67 \\
low & 0.20 & 0.112 & 0.30 & 0.151 & 0.03 & 0.32 & 0.65 \\
trusted & 0.01 & 0.005 & 0.27 & 0.005 & 0.03 & 0.63 & 0.67 \\
trusted & 0.05 & 0.027 & 0.28 & 0.033 & 0.03 & 0.66 & 0.66 \\
trusted & 0.20 & 0.115 & 0.30 & 0.153 & 0.03 & 0.63 & 0.63 \\
\bottomrule
\end{tabular}
\end{table}
Table~\ref{tab:poison} gives the secondary result. At 20\% injection through a low-trust channel, \pal{} answers \emph{established} with the poisoned value on \PoisHarmPal\% of value queries against \PoisHarmTdb\% for the temporal DB and \PoisHarmCommit\% for always-commit; the difference is smaller than one might expect from a justification-based store because a poisoned report carrying a false change cue is shielded from earlier competitors under A-CUE, so cue resolution is an attack surface. The support layer discriminates poisoned claims through the low channel (mean support \PoisSupLowPois{} vs.\ \PoisSupLowHon{} for honest candidates) but not through a compromised trusted channel (\PoisSupTrPois{} vs.\ \PoisSupTrHon): provenance cannot detect what provenance has endorsed, and only corroboration across origins could. Both effects are what ``measured, not assumed from provenance'' was meant to expose.

\subsection{Setting 2: natural language, automatic extraction}
\begin{sidewaystable}\centering\footnotesize\setlength{\tabcolsep}{3.5pt}
\caption{Setting 2 (natural-language streams, automatic extraction with Claude Haiku 4.5; 100 streams, 99 with a warranted downstream query). Stream-level means with 95\% cluster-bootstrap intervals; cost is measured tokens (input $+8\times$ output) per stream.}\label{tab:s2}
\begin{tabular}{lccccccc}
\toprule
System & Warr.\ downstream & Warr.\ direct & Resolvable & Unresolved & Commit rate$\downarrow$ & Prov.\ P & Cost (tok/stream) \\
\midrule
\textsc{Palimpsest} & 0.956 {\tiny[0.93,0.98]} & 0.996 {\tiny[0.99,1.00]} & 0.976 {\tiny[0.97,0.98]} & 0.834 {\tiny[0.80,0.87]} & 0.080 {\tiny[0.06,0.10]} & 0.978 {\tiny[0.97,0.98]} & 188k \\
\quad no propagation & 0.749 {\tiny[0.70,0.79]} & 0.996 {\tiny[0.99,1.00]} & 0.920 {\tiny[0.91,0.93]} & 0.813 {\tiny[0.78,0.85]} & 0.079 {\tiny[0.06,0.10]} & 0.972 {\tiny[0.97,0.98]} & 188k \\
\quad forced commit & 0.979 {\tiny[0.96,0.99]} & 0.996 {\tiny[0.99,1.00]} & 0.987 {\tiny[0.98,0.99]} & 0.045 {\tiny[0.03,0.06]} & 0.758 {\tiny[0.73,0.79]} & 0.985 {\tiny[0.98,0.99]} & 188k \\
Symbolic replay (rule adjudicator) & 0.956 {\tiny[0.93,0.98]} & 0.996 {\tiny[0.99,1.00]} & 0.976 {\tiny[0.97,0.98]} & 0.834 {\tiny[0.80,0.87]} & 0.080 {\tiny[0.06,0.10]} & 0.989 {\tiny[0.99,0.99]} & 188k \\
Claims log + LLM adjudicator & 0.514 {\tiny[0.46,0.56]} & 0.731 {\tiny[0.69,0.77]} & 0.666 {\tiny[0.64,0.69]} & 0.128 {\tiny[0.11,0.15]} & 0.572 {\tiny[0.53,0.61]} & 0.747 {\tiny[0.73,0.76]} & 532k \\
\quad top-40 entries (original, 30 str.) & 0.419 {\tiny[0.32,0.52]} & 0.731 {\tiny[0.64,0.81]} & 0.614 {\tiny[0.57,0.66]} & 0.103 {\tiny[0.08,0.13]} & 0.504 {\tiny[0.46,0.54]} & 0.758 {\tiny[0.73,0.78]} & 454k \\
Raw-report log + LLM adjudicator (30 str.) & 0.330 {\tiny[0.24,0.43]} & 0.685 {\tiny[0.59,0.78]} & 0.554 {\tiny[0.51,0.60]} & 0.101 {\tiny[0.07,0.13]} & 0.606 {\tiny[0.55,0.66]} & 0.683 {\tiny[0.66,0.71]} & 439k \\
Claims log + LLM adj.\ (amended contract, 30 str.) & 0.400 {\tiny[0.33,0.47]} & 0.703 {\tiny[0.60,0.80]} & 0.607 {\tiny[0.57,0.65]} & 0.134 {\tiny[0.09,0.18]} & 0.533 {\tiny[0.47,0.59]} & 0.714 {\tiny[0.69,0.74]} & 552k \\
\quad top-10 entries & 0.225 {\tiny[0.16,0.29]} & 0.787 {\tiny[0.71,0.86]} & 0.572 {\tiny[0.53,0.62]} & 0.138 {\tiny[0.10,0.17]} & 0.377 {\tiny[0.33,0.43]} & 0.836 {\tiny[0.82,0.85]} & 399k \\
\quad top-20 entries & 0.385 {\tiny[0.30,0.48]} & 0.723 {\tiny[0.63,0.81]} & 0.610 {\tiny[0.57,0.65]} & 0.136 {\tiny[0.10,0.18]} & 0.412 {\tiny[0.36,0.46]} & 0.786 {\tiny[0.77,0.80]} & 424k \\
\quad top-40 entries & 0.429 {\tiny[0.33,0.53]} & 0.739 {\tiny[0.65,0.82]} & 0.628 {\tiny[0.59,0.67]} & 0.133 {\tiny[0.10,0.17]} & 0.471 {\tiny[0.42,0.52]} & 0.748 {\tiny[0.73,0.77]} & 464k \\
Raw-report log + LLM adj.\ (amended, 30 str.) & 0.323 {\tiny[0.23,0.42]} & 0.672 {\tiny[0.59,0.76]} & 0.524 {\tiny[0.48,0.57]} & 0.102 {\tiny[0.07,0.14]} & 0.563 {\tiny[0.51,0.61]} & 0.674 {\tiny[0.65,0.70]} & 448k \\
Temporal DB & 0.779 {\tiny[0.74,0.82]} & 0.885 {\tiny[0.85,0.92]} & 0.895 {\tiny[0.88,0.91]} & 0.000 {\tiny[0.00,0.00]} & 0.837 {\tiny[0.81,0.87]} & 0.961 {\tiny[0.95,0.97]} & 188k \\
Always abstain & 0.000 {\tiny[0.00,0.00]} & 0.000 {\tiny[0.00,0.00]} & 0.000 {\tiny[0.00,0.00]} & 0.380 {\tiny[0.35,0.41]} & 0.000 {\tiny[0.00,0.00]} & 0.542 {\tiny[0.52,0.56]} & 188k \\
Always commit & 0.422 {\tiny[0.38,0.46]} & 0.982 {\tiny[0.97,0.99]} & 0.791 {\tiny[0.77,0.81]} & 0.000 {\tiny[0.00,0.00]} & 0.795 {\tiny[0.77,0.82]} & 0.947 {\tiny[0.94,0.96]} & 188k \\
\bottomrule
\end{tabular}
\end{sidewaystable}
\paragraph{Evaluation set.} 100 streams (13{,}094 reports) were rendered, paraphrased by Claude Haiku 4.5 and passed through the cue-recovery check; after two regenerations 56 lines had lost a stated change and 3 had gained one, and 1{,}343 were flagged by the model-side check as reading a dated statement (``since 3 March'') as a change---flags that cannot alter gold but are counted. Adjudication stripped the lost cues from the frozen labels, changing 274 of 6{,}737 gold entries (4.1\%); the frozen set has \sTwoWarrPerStream{} warranted downstream queries per stream, just under the pre-registered floor of 8 (deviation recorded). The shared extractor (Haiku 4.5) has a slot-error rate of \sTwoExtSlot\%, misses a stated cue on \sTwoExtMissed\% of lines, adds none (\sTwoExtSpur\%), errs on dates on \sTwoExtValid\%, and fails on \sTwoExtFail\%.

\paragraph{H1.} On warranted downstream queries, \pal{} over the extracted claims reaches \sTwoPalWarr{} against \sTwoLogWarr{} for the same claims log adjudicated by the LLM at query time: a paired difference of $+\sTwoHOne$ points $[\sTwoHOneLo, \sTwoHOneHi]$ over \sTwoHOneN{} streams, far beyond the pre-registered 10-point threshold, with the interval excluding zero. The gap is not an artefact of the structured schema: the rule-based adjudicator over the same claims (the upper bound of what query-time reasoning could achieve) reaches \sTwoRuleWarr{}, so the LLM adjudicator loses 44 points to reasoning, not to extraction. Removing propagation costs $\sTwoProp$ points $[\sTwoPropLo, \sTwoPropHi]$, the same effect as in Setting 1. The raw-report log with no extraction at all, the pure query-time alternative, reaches \sTwoNlWarr{} on its 30 streams, and the matched-cost variant that shows the adjudicator only the 40 most recent entries plus the queried key reaches \sTwoTopkWarr{}: less context makes the adjudicator worse, not better.

\paragraph{H4.} The LLM adjudicator commits on \sTwoLogCommit{} of unresolved slots against \sTwoPalCommit{} for \pal, while \pal's resolvable accuracy is \sTwoPalRes{} against \sTwoLogRes: both halves of the joint criterion hold. The adjudicator also changes downstream answers after non-warranted events \sTwoLogUnsup{} of the time (\pal: \sTwoPalUnsup). The forced-commit variant is instructive: on extracted claims it scores \sTwoCommitWarr{} on warranted queries, slightly \emph{above} \pal{} (\sTwoPalWarr), because committing to the minimal-error interpretation often resolves a spurious dispute created by an extraction error in the right direction---but it does so by committing on \sTwoCommitCommit{} of genuinely unresolved slots, which is the trade the contract makes visible.

\paragraph{Cost and the crossover (H2).} Costs are measured tokens. Extraction costs \xE{}k tokens per stream (\$\sTwoPalCost{} at batch prices) and \pal{} pays nothing per query; the claims-log adjudicator pays \xAlog{}k tokens per query on top (\$\sTwoLogCost{} per stream at this query load), the raw-report adjudicator \xAnl{}k per query with no extraction. The crossover ratio at which write-time extraction becomes cheaper than query-time adjudication over raw reports is $r^*=\xRstar$ queries per observation; below it the raw log is cheaper, but at every ratio it is far less accurate, so the accuracy--cost curves do not cross in accuracy. With a shared extractor the claims log is strictly dearer than \pal{} at any query load, and less accurate.

\paragraph{Sign stability with a stronger model.}
\begin{sidewaystable}\centering\footnotesize\setlength{\tabcolsep}{3.5pt}
\caption{Sign-stability subset: 15 Setting 2 streams re-extracted and adjudicated with Claude Sonnet 4.6 on the same frozen evaluation set.}\label{tab:s2strong}
\begin{tabular}{lccccccc}
\toprule
System & Warr.\ downstream & Warr.\ direct & Resolvable & Unresolved & Commit rate$\downarrow$ & Prov.\ P & Cost (tok/stream) \\
\midrule
\textsc{Palimpsest} & 1.000 {\tiny[1.00,1.00]} & 1.000 {\tiny[1.00,1.00]} & 1.000 {\tiny[1.00,1.00]} & 0.816 {\tiny[0.73,0.90]} & 0.090 {\tiny[0.05,0.14]} & 0.987 {\tiny[0.98,1.00]} & 189k \\
\quad no propagation & 0.780 {\tiny[0.64,0.90]} & 1.000 {\tiny[1.00,1.00]} & 0.946 {\tiny[0.91,0.98]} & 0.805 {\tiny[0.72,0.89]} & 0.090 {\tiny[0.05,0.14]} & 0.984 {\tiny[0.97,0.99]} & 189k \\
\quad forced commit & 1.000 {\tiny[1.00,1.00]} & 1.000 {\tiny[1.00,1.00]} & 1.000 {\tiny[1.00,1.00]} & 0.048 {\tiny[0.00,0.12]} & 0.751 {\tiny[0.62,0.87]} & 0.989 {\tiny[0.98,1.00]} & 189k \\
Symbolic replay (rule adjudicator) & 1.000 {\tiny[1.00,1.00]} & 1.000 {\tiny[1.00,1.00]} & 1.000 {\tiny[1.00,1.00]} & 0.816 {\tiny[0.73,0.90]} & 0.090 {\tiny[0.05,0.14]} & 0.997 {\tiny[0.99,1.00]} & 189k \\
Claims log + LLM adjudicator & 0.813 {\tiny[0.67,0.93]} & 0.924 {\tiny[0.85,0.98]} & 0.854 {\tiny[0.80,0.91]} & 0.289 {\tiny[0.19,0.39]} & 0.489 {\tiny[0.36,0.63]} & 0.851 {\tiny[0.82,0.88]} & 718k \\
Temporal DB & 0.777 {\tiny[0.64,0.89]} & 0.844 {\tiny[0.67,0.97]} & 0.908 {\tiny[0.86,0.95]} & 0.000 {\tiny[0.00,0.00]} & 0.833 {\tiny[0.72,0.92]} & 0.957 {\tiny[0.94,0.97]} & 189k \\
Always abstain & 0.000 {\tiny[0.00,0.00]} & 0.000 {\tiny[0.00,0.00]} & 0.000 {\tiny[0.00,0.00]} & 0.399 {\tiny[0.30,0.49]} & 0.000 {\tiny[0.00,0.00]} & 0.531 {\tiny[0.47,0.58]} & 189k \\
Always commit & 0.387 {\tiny[0.29,0.48]} & 0.990 {\tiny[0.97,1.00]} & 0.778 {\tiny[0.72,0.83]} & 0.000 {\tiny[0.00,0.00]} & 0.879 {\tiny[0.84,0.92]} & 0.945 {\tiny[0.92,0.97]} & 189k \\
\bottomrule
\end{tabular}
\end{sidewaystable}
Re-extracting and adjudicating 15 streams with Claude Sonnet 4.6 on the same frozen evaluation set (Table~\ref{tab:s2strong}) lowers the extraction slot-error rate to \sStrongExtSlot\% and lifts the adjudicator to \sStrongLogWarr{} on warranted downstream queries; \pal{} reaches \sStrongPalWarr{} and the paired difference remains positive, $+\sStrongHOne$ points $[\sStrongHOneLo, \sStrongHOneHi]$, with the adjudicator still committing on \sStrongLogCommit{} of unresolved slots. The sign is stable; the size of the gap depends strongly on the adjudicator.

\paragraph{RQ2: extraction noise.}
\begin{table}[t]\centering\small
\caption{RQ2 corruption curve (50 Setting 2 streams, mock extractor with injected error rates). Warranted downstream accuracy; PALIMPSEST resolvable and unresolved accuracy.}\label{tab:noise}
\begin{tabular}{llcccccc}
\toprule
Error type & rate & \textsc{Palimpsest} & no prop. & Temporal DB & Always commit & PAL resolvable & PAL unresolved \\
\midrule
wrong value & 0.00 & 1.000 & 0.801 & 0.799 & 0.423 & 1.000 & 1.000 \\
wrong value & 0.05 & 0.878 & 0.706 & 0.738 & 0.404 & 0.910 & 0.858 \\
wrong value & 0.10 & 0.658 & 0.519 & 0.625 & 0.354 & 0.814 & 0.685 \\
wrong value & 0.20 & 0.449 & 0.332 & 0.388 & 0.257 & 0.652 & 0.515 \\
wrong value & 0.35 & 0.273 & 0.192 & 0.301 & 0.224 & 0.511 & 0.347 \\
change cue dropped & 0.10 & 0.912 & 0.713 & 0.797 & 0.423 & 0.946 & 0.958 \\
change cue dropped & 0.25 & 0.817 & 0.624 & 0.794 & 0.423 & 0.867 & 0.873 \\
change cue dropped & 0.50 & 0.562 & 0.382 & 0.758 & 0.423 & 0.712 & 0.753 \\
change cue added & 0.05 & 1.000 & 0.801 & 0.799 & 0.423 & 1.000 & 0.963 \\
change cue added & 0.15 & 1.000 & 0.815 & 0.799 & 0.423 & 1.000 & 0.871 \\
date dropped & 0.25 & 0.995 & 0.798 & 0.794 & 0.423 & 0.995 & 0.973 \\
date dropped & 0.50 & 0.991 & 0.794 & 0.790 & 0.423 & 0.990 & 0.961 \\
\bottomrule
\end{tabular}
\end{table}
Table~\ref{tab:noise} injects extraction errors at controlled rates with a mock extractor on 50 streams. \pal{} degrades faster than last-write-wins as wrong values accumulate: at 20\% wrong values it scores \nzSlotTwentyPal{} against \nzSlotTwentyTdb{} for the temporal DB, and at 35\% the order reverses (\nzSlotThirtyfivePal{} vs.\ \nzSlotThirtyfiveTdb); when half the stated change cues are dropped, the temporal DB (\nzCueFiftyTdb) is clearly ahead of \pal{} (\nzCueFiftyPal). The mechanism is simple and, we think, the most important limitation of justified-belief maintenance: every extraction error becomes a standing dispute that the semantics keeps open, whereas last-write-wins silently discards it. Dropped dates and added cues cost almost nothing. At the measured error rates of the real extractors (0.3--1.5\%) the effect is small, but it bounds the regime in which the approach pays.

\paragraph{H3: calibration on extracted claims.} For the \RqTwoSupN{} value queries the LLM adjudicator answered as established, the evidence model's support for the asserted value (computed from the extracted evidence) has ECE \RqTwoSupEce{} and Brier \RqTwoSupBrier, worse than the adjudicator's own verbalised confidence (ECE \RqTwoVerbEce, Brier \RqTwoVerbBrier) and than agreement across five samples at temperature 0.7 on a 300-query subset (ECE \RqTwoSampEce, Brier \RqTwoSampBrier); only the risk--coverage areas are close (\RqTwoSupAurc, \RqTwoVerbAurc, \RqTwoSampAurc). H3 is therefore \emph{not} supported in Setting 2: on this population the model's confidence is better calibrated than a provenance-and-recency rule. Together with the Setting 1 comparison against a learned predictor, the evidence model's claim reduces to a defined, auditable meaning of support, not to calibration.

\subsection{Setting 3: streams assembled from Wikidata statement histories}
\begin{sidewaystable}\centering\footnotesize\setlength{\tabcolsep}{3.5pt}
\caption{Setting 3 (transfer: 30 streams assembled from Wikidata statement histories, pseudonymised surface, extraction and adjudication with Claude Haiku 4.5). Descriptive; the warranted-downstream stratum has 24 queries in 13 streams.}\label{tab:s3}
\begin{tabular}{lccccccc}
\toprule
System & Warr.\ downstream & Warr.\ direct & Resolvable & Unresolved & Commit rate$\downarrow$ & Prov.\ P & Cost (tok/stream) \\
\midrule
\textsc{Palimpsest} & 0.923 {\tiny[0.77,1.00]} & 0.976 {\tiny[0.94,1.00]} & 0.978 {\tiny[0.95,1.00]} & 0.909 {\tiny[0.73,1.00]} & 0.000 {\tiny[0.00,0.00]} & 0.892 {\tiny[0.85,0.94]} & 25k \\
\quad no propagation & 0.724 {\tiny[0.49,0.93]} & 0.976 {\tiny[0.94,1.00]} & 0.962 {\tiny[0.93,0.99]} & 0.879 {\tiny[0.70,1.00]} & 0.000 {\tiny[0.00,0.00]} & 0.886 {\tiny[0.84,0.93]} & 25k \\
\quad forced commit & 0.923 {\tiny[0.77,1.00]} & 0.976 {\tiny[0.94,1.00]} & 0.978 {\tiny[0.95,1.00]} & 0.000 {\tiny[0.00,0.00]} & 1.000 {\tiny[1.00,1.00]} & 0.892 {\tiny[0.85,0.94]} & 25k \\
Symbolic replay (rule adjudicator) & 0.923 {\tiny[0.77,1.00]} & 0.976 {\tiny[0.94,1.00]} & 0.978 {\tiny[0.95,1.00]} & 0.909 {\tiny[0.73,1.00]} & 0.000 {\tiny[0.00,0.00]} & 1.000 {\tiny[1.00,1.00]} & 25k \\
Claims log + LLM adjudicator & 0.019 {\tiny[0.00,0.06]} & 0.519 {\tiny[0.38,0.66]} & 0.185 {\tiny[0.14,0.23]} & 0.333 {\tiny[0.09,0.61]} & 0.545 {\tiny[0.27,0.79]} & 0.802 {\tiny[0.75,0.86]} & 88k \\
Raw-report log + LLM adjudicator (30 str.) & 0.077 {\tiny[0.00,0.23]} & 0.451 {\tiny[0.30,0.61]} & 0.181 {\tiny[0.13,0.23]} & 0.182 {\tiny[0.00,0.45]} & 0.576 {\tiny[0.30,0.82]} & 0.731 {\tiny[0.67,0.79]} & 71k \\
Claims log + LLM adj.\ (amended contract, 30 str.) & 0.019 {\tiny[0.00,0.06]} & 0.491 {\tiny[0.34,0.64]} & 0.205 {\tiny[0.15,0.26]} & 0.394 {\tiny[0.15,0.64]} & 0.485 {\tiny[0.21,0.73]} & 0.791 {\tiny[0.73,0.85]} & 92k \\
Raw-report log + LLM adj.\ (amended, 30 str.) & 0.173 {\tiny[0.00,0.38]} & 0.466 {\tiny[0.32,0.61]} & 0.208 {\tiny[0.15,0.26]} & 0.318 {\tiny[0.09,0.59]} & 0.439 {\tiny[0.18,0.70]} & 0.728 {\tiny[0.67,0.79]} & 25k \\
Temporal DB & 0.865 {\tiny[0.65,1.00]} & 0.884 {\tiny[0.78,0.97]} & 0.692 {\tiny[0.64,0.74]} & 0.000 {\tiny[0.00,0.00]} & 1.000 {\tiny[1.00,1.00]} & 0.922 {\tiny[0.88,0.96]} & 25k \\
Always abstain & 0.000 {\tiny[0.00,0.00]} & 0.000 {\tiny[0.00,0.00]} & 0.000 {\tiny[0.00,0.00]} & 0.788 {\tiny[0.55,1.00]} & 0.000 {\tiny[0.00,0.00]} & 0.828 {\tiny[0.78,0.88]} & 25k \\
Always commit & 0.000 {\tiny[0.00,0.00]} & 0.506 {\tiny[0.36,0.65]} & 0.478 {\tiny[0.43,0.53]} & 0.000 {\tiny[0.00,0.00]} & 0.909 {\tiny[0.73,1.00]} & 0.803 {\tiny[0.75,0.86]} & 25k \\
\bottomrule
\end{tabular}
\end{sidewaystable}
Setting 3 is the transfer setting: 30 streams that nobody on the project designed, assembled from the revision histories of 40 Wikidata persons and the 128 organisations recorded as their employers (3,474 sampled revisions; selection, mapping and schema decisions are listed in the released note \texttt{docs/SETTING3.md}). Every edit between sampled revisions becomes an observation: a claim added is an assertion (with a \emph{since} cue when it carries a start qualifier), a value edited in place is a correction, a statement deprecated or removed by the editors is a withdrawal. References give the source class (unreferenced statements are \emph{low}; 331 of 417 assertions). The surface is pseudonymised with the same fictional names as Settings 1--2, so the backbone cannot answer from knowledge of the real person; the timing, conflicts and withdrawals are real. One schema decision matters and is ours: Wikidata's employer property lists every affiliation over a career, mostly undated, so it is declared a stable, multi-valued set (under a single-valued or a changeable reading the derived slot is unresolved almost everywhere: 1 warranted downstream query in 30 streams). The derived slot is therefore the set of headquarters cities of the recorded affiliations, and a warranted downstream query is one where a withdrawal or correction of an affiliation, or a headquarters change, alters that set. The resulting set has 497 observations, 80 withdrawals, \sThreeWarrTotal{} warranted downstream queries in 13 streams (\sThreeWarrPerStream{} per stream), and is reported descriptively, as pre-registered.

Table~\ref{tab:s3} gives the same systems and pipeline as Setting 2 (Haiku 4.5 extraction, slot-error rate \sThreeExtSlot\%). \pal{} over the extracted claims reaches \sThreePalWarr{} on warranted downstream queries, equal to the rule-based adjudicator over the same claims (\sThreeRuleWarr): the ceiling set by extraction. Removing propagation costs $\sThreeProp$ points $[\sThreePropLo, \sThreePropHi]$ over 13 streams, so retraction propagation matters on real histories as it did on generated ones. The LLM adjudicator reaches \sThreeLogWarr{} on the same queries and \sThreeLogRes{} on resolvable ones, a paired difference of $+\sThreeHOne$ points $[\sThreeHOneLo, \sThreeHOneHi]$, and we decompose this gap because most of it is not reasoning in the sense Setting 2 measured. Two output conventions account for a large part: \sThreeEstEmpty{} of the \sThreeEstTotal{} established queries have the empty set as gold (a multi-valued slot whose every reading justifies no member, typically after the last affiliation is withdrawn), which the contract as run did not state and on which the adjudicator always answered \emph{unknown}; and \sThreeHoldsFalse{} yes/no queries are established-false (a value that never held at the asked time, frequent with \emph{since} cues), which the contract does state but the adjudicator again answered \emph{unknown}. Excluding both, its resolvable accuracy is \sThreeLogResNe{} (\pal: \sThreePalResNe). On the warranted downstream queries themselves it answers \emph{unresolved} with the gold set among its alternatives in most cases: it reasons that overlapping affiliations ``may have ended'' although the declared schema marks the attribute stable, i.e.\ it applies its prior about employment to a schema that contradicts it. The temporal DB, which has no such prior, reaches \sThreeTdbWarr{} on warranted downstream queries and commits on every unresolved slot. We draw three conclusions: write-time maintenance transfers to real revision histories (the propagation effect and the equality with the rule adjudicator); the output contract must state the empty-set convention, which we have now added (dated); and the LLM adjudicator's failure here is partly a schema-versus-prior conflict that a pseudonymised, declared schema makes unusually sharp---a limitation of this transfer design as much as of the adjudicator.

\subsection{External check: LongMemEval knowledge updates}
\begin{table}[t]\centering\footnotesize
\caption{LongMemEval knowledge-update items (78; 69 mapped to an extracted attribute). Systems answer one structured slot from claims extracted turn by turn; a Haiku judge compares the answer with the reference. Coverage = share answered as established; accuracy over all items and over answered items. Not comparable to published LongMemEval scores.}\label{tab:lme}
\begin{tabular}{lccc}
\toprule
System & Coverage & Acc.\ (all) & Acc.\ (answered) \\
\midrule
\pal{} (P0c) & 0.32 & 0.167 & 0.400 \\
\pal{} + self-update (exploratory) & 0.74 & 0.474 & 0.621 \\
Temporal DB (last write wins) & 0.88 & 0.526 & 0.594 \\
Always-commit & 0.88 & 0.526 & 0.594 \\
Always-abstain & 0.00 & 0.038 & 0.000 \\
Claims log + LLM adjudicator & 0.56 & 0.359 & 0.614 \\
\bottomrule
\end{tabular}
\end{table}
LongMemEval's knowledge-update items \citep{longmemeval} ask for a user attribute that was stated and later updated across chat sessions. We run the \lmeN{} items of the oracle release through an open-vocabulary adapter shared by all systems: each user turn is extracted into claims about the user (Haiku), attribute keys are canonicalised per item, the question is mapped to one key, and every system answers that slot; a Haiku judge compares the answer with the reference, and on the six items whose reference is ``not enough information'' abstention is credited. These numbers are not comparable to published LongMemEval scores (the systems answer a structured slot, not the question), and the adapter caps all systems: 9 items map to no extracted key.

The result (Table~\ref{tab:lme}) is a scoring-target mismatch rather than a measurement of memory quality, and we report it as such. The benchmark's reference for an update is the user's \emph{latest} statement; nothing in the data verifies it, and the extractor recovered a stated change on 5 of 91 claims on the queried keys. Under our semantics a later uncued value from the same source is a dispute, so \pal{} answers \emph{unresolved} on \lmePalUnres{} of 69 mapped items (coverage \lmePalCov, accuracy \lmePalAcc), where last-write-wins (the temporal DB) answers and is right on 34 of them (accuracy \lmeTdbAcc, coverage \lmeTdbCov); on the items where both commit they agree. The LLM adjudicator abstains on 44\% and reaches \lmeLogAcc. An exploratory variant that reads a later self-report on a single-valued key as a stated change (added after seeing this result, not pre-registered, shown only to make the mechanism visible) recovers \lmePalSuAcc{} at coverage \lmePalSuCov. The benchmark therefore measures the recency convention for single-source self-reports, which is a reasonable convention for that setting and which our semantics deliberately does not assume; the ten items where recency committed and was judged wrong are mostly questions about previous values or about the change itself, which a single current-value slot cannot express, and three presupposition-false questions where recency answers anyway. What LongMemEval does not contain---unreliable sources, hidden world state, withdrawn evidence---is what Settings 1--3 measure. STALE has no public data release, so no further external check was possible.


\section{Threats to validity}\label{sec:threats}
\paragraph{Construct.} The generator, the semantics and the system share an author; Settings 1--2 could be tuned to the method. The hand-written streams, the declared perturbation menu, the pre-registered selection rule and the sensitivity variants are the defence we can offer; the transfer setting and the external check are the only independent ones, and they are small. The transfer setting is independent in its data, not in its mapping: the selection query asked for histories with deprecated statements and ended employments, and the schema (notably the stable affiliation set) is ours; a different mapping would give a different stratum, and the one we chose was reached after two mappings left the derived slot unresolved almost everywhere, which we report. The semantics itself embodies choices a reviewer may reject: inertia (A4), report-time anchoring (A1), and above all cue resolution (A-CUE), which we adopted after measuring that literal P0 leaves the benchmark abstention-dominated, and which the poisoning experiment shows to be exploitable.
\paragraph{Internal.} The shared extractor and identical instruction remove the most obvious asymmetry between \pal{} and the log baseline, but a prompt can still suit a structured schema better than free adjudication; we report the rule-adjudicated log as the upper bound of what the baseline's reasoning could achieve with the same claims. The provenance metric assumes the oracle's justification paths are the right notion of support.
\paragraph{Statistical.} 500 and 150 streams give narrow intervals on the primary endpoints; the transfer setting (30 streams) and the external subsets are reported as evidence, not tested. All intervals are cluster bootstraps at the stream level because queries within a stream share errors.
\paragraph{External.} One backbone family at full scale; the robustness subset tests sign stability only. Fictional entity names remove world knowledge from the natural-language settings, which is a deliberate control and a limitation: real reports lean on it, and in Setting 3 the pseudonymised schema put the adjudicator's priors about employment against a declaration it was told to follow. The output contract as run did not state the empty-set convention for multi-valued slots; it is stated now (dated), and the affected queries (\sThreeEstEmpty{} in Setting 3, 64 in Setting 2) are counted against the LLM adjudicator in the tables and disclosed in the text. LongMemEval scores recency on single-source self-reports, which our semantics does not assume; the gap there is a mismatch of scoring targets, and it also shows that the extractor recovers stated changes poorly when they are implicit. Negative-polarity assertions and interval valid-time cues are deferred.
\paragraph{Budget-induced.} The whole study ran under a \$100 ceiling; STALE could not be used at its native context lengths and was not released as data in any case; the external checks are subsamples whose numbers are not comparable to published headline figures.


\section{Conclusion}
We set out to answer a narrow question with instruments that cannot reward the easy failure modes: a semantics of justified belief that keeps alternatives open, a contract that scores behaviour rather than representation, and a benchmark that separates what happened from what was reported and from what the reports justify. In the structured setting the instruments discriminate as designed, retraction propagation is isolated from ordinary updating, and two honest negatives emerged: a learned predictor calibrates better than the hand-designed evidence model, and cue resolution---needed to keep the benchmark from being abstention-dominated---is an attack surface that provenance does not close. In the natural-language setting the primary hypothesis held by a wide margin and at lower cost, the matched-cost and raw-report variants of query-time adjudication did worse rather than better, and the sign survived a stronger adjudicator; the calibration hypothesis failed, and the corruption curve located the regime where last-write-wins overtakes justified belief. On streams assembled from real revision histories the propagation effect and the equality with a rule-based adjudicator transfer; the LLM adjudicator's much larger gap there is in good part a conflict between its priors and a declared schema, and a convention the contract failed to state, both of which we report rather than net out. On LongMemEval, a benchmark that scores recency on self-reports, a semantics that refuses to assume recency is, as it should be, penalised; what that benchmark does not contain is what this study is about. Everything needed to reproduce, contest or extend these results, including the pre-registration with its dated deviations and the cost ledger, is released.


\bibliographystyle{plainnat}
\bibliography{refs}
\end{document}
